\documentclass[12pt,oneside]{book}

% -------------------------------------------------
% PAQUETES
% -------------------------------------------------
\usepackage[spanish,es-nodecimaldot]{babel}
\usepackage[utf8]{inputenc}
\usepackage[T1]{fontenc}
\usepackage{lmodern}
\usepackage[a4paper,margin=2.5cm]{geometry}
\usepackage{microtype}
\usepackage{setspace}
\usepackage{graphicx}
\usepackage{xcolor}
\usepackage{fancyhdr}
\usepackage{titlesec}
\usepackage{hyperref}
\usepackage{bookmark}
\usepackage{tocloft}
\usepackage{tcolorbox}
\usepackage{amsmath,amssymb}
\usepackage{physics}
\usepackage{enumitem}
\usepackage{csquotes}
\usepackage{listings}

% -------------------------------------------------
% CONFIGURACIÓN VISUAL
% -------------------------------------------------
\definecolor{verdeUdeA}{HTML}{007A33}
\definecolor{grisClaro}{HTML}{F4F6F8}
\definecolor{grisTexto}{HTML}{2E2E2E}
\definecolor{grisSec}{HTML}{6C757D}

\hypersetup{
    colorlinks=true,
    linkcolor=verdeUdeA,
    urlcolor=verdeUdeA,
    citecolor=verdeUdeA,
    pdftitle={Notas de Clase - Mecánica Celeste},
    pdfauthor={Daniel Soto Villada}
}

\setstretch{1.12}
\setlength{\parskip}{0.4em}
\setlength{\parindent}{0pt}
\setcounter{secnumdepth}{3}
\setcounter{tocdepth}{2}
\setlength{\headheight}{38pt}

% Estilo de capítulos y secciones
\titleformat{\chapter}[display]
  {\bfseries\Huge\color{verdeUdeA}}
  {\filleft\Large\chaptername\ \thechapter}
  {1ex}
  {\titlerule\vspace{1ex}\filright}
  [\vspace{1ex}\titlerule]

\titleformat{\section}
  {\bfseries\Large\color{verdeUdeA}}
  {\thesection}{0.7em}{}

\titleformat{\subsection}
  {\bfseries\large\color{grisTexto}}
  {\thesubsection}{0.6em}{}

% Encabezados y pies
\pagestyle{fancy}
\fancyhf{}
\fancyhead[L]{\textcolor{grisSec}{\nouppercase{\leftmark}}}
\fancyhead[R]{\textcolor{grisSec}{Daniel Soto Villada}}
\fancyfoot[C]{\textcolor{grisSec}{\thepage}}
\renewcommand{\headrulewidth}{0.4pt}
\renewcommand{\headrule}{\hbox to\headwidth{\color{verdeUdeA}\leaders\hrule height \headrulewidth\hfill}}

% Índice más limpio
\renewcommand{\cftchapfont}{\bfseries\color{verdeUdeA}}
\renewcommand{\cftsecfont}{\color{grisTexto}}
\renewcommand{\cftchappagefont}{\bfseries\color{verdeUdeA}}
\renewcommand{\cftsecpagefont}{\color{grisTexto}}

% Caja para definiciones/ideas importantes
\newtcolorbox{importante}{
  colback=grisClaro,
  colframe=verdeUdeA,
  boxrule=0.8pt,
  arc=2mm,
  left=2mm,right=2mm,top=1mm,bottom=1mm
}

% -------------------------------------------------
% DATOS DEL DOCUMENTO
% -------------------------------------------------
\newcommand{\NombreEstudiante}{Daniel Soto Villada}
\newcommand{\Programa}{Estudiante de astronomía}
\newcommand{\Universidad}{Universidad de Antioquia}
\newcommand{\Facultad}{Facultad de Ciencias Exactas y Naturales}
\newcommand{\Instituto}{Instituto de Física}
\newcommand{\TituloDocumento}{Notas de Clase - Mecánica Celeste}

% -------------------------------------------------
% DOCUMENTO
% -------------------------------------------------
\begin{document}
\hypersetup{pageanchor=false}

% ---------- PORTADA ----------
\begin{titlepage}
    \centering
    \vspace*{1.8cm}

    {\Large\textcolor{grisSec}{\Universidad}\par}
    \vspace{0.25cm}
    {\large\textcolor{grisSec}{\Facultad}\par}
    {\large\textcolor{grisSec}{\Instituto}\par}

    \vspace{2.2cm}
    {\Huge\bfseries\textcolor{verdeUdeA}{\TituloDocumento}\par}


    \vspace{2.4cm}
    \begin{tcolorbox}[colback=grisClaro,colframe=verdeUdeA,boxrule=0.9pt,arc=2mm,width=0.78\textwidth]
        \centering
        {\Large\bfseries \NombreEstudiante}\par
        \vspace{0.2cm}
        {\large \Programa}
    \end{tcolorbox}

    \vfill
    {\large\textcolor{grisSec}{\today}\par}
    \vspace{0.5cm}
\end{titlepage}

% ---------- PÁGINAS PRELIMINARES ----------
\frontmatter
\tableofcontents
\clearpage

% ---------- CONTENIDO ----------
\clearpage
\hypersetup{pageanchor=true}
\mainmatter
\input{clases/introduccion}
\input{clases/clase13}
\input{clases/clase14}
\input{clases/clase15}
\chapter{Clase 16: Cuadraturas del problema relativo de dos cuerpos}

Esta clase continúa el desarrollo analítico del problema relativo de dos cuerpos introducido en la clase anterior. El objetivo es identificar las \textbf{integrales primeras} o \textbf{cuadraturas} que permiten resolver completamente la dinámica relativa y establecer por qué este problema es integrable de forma exacta.

\section{Planteamiento del sistema y necesidad de cuadraturas}
Partimos de la ecuación diferencial que describe el movimiento relativo entre dos cuerpos sometidos a gravitación mutua:
\[
\ddot{\vec{r}} = -\frac{\mu}{r^3}\vec{r}
\]
donde \(\mu = G(m_1 + m_2)\) y \(r = |\vec{r}|\).

Esta es una ecuación diferencial vectorial de segundo orden. Al descomponerla en coordenadas cartesianas, el problema queda gobernado por \textbf{6 variables dinámicas independientes} \((x, y, z, \dot{x}, \dot{y}, \dot{z})\). Para integrar completamente el sistema y obtener una solución cerrada, es necesario hallar \textbf{6 cuadraturas} o integrales primeras independientes.

Aunque el problema de \(N\) cuerpos posee constantes de movimiento globales, el problema relativo de dos cuerpos constituye un sistema autónomo nuevo que requiere su propia deducción analítica paso a paso.

\begin{importante}
La meta de esta clase no es todavía escribir la órbita final en forma cónica, sino mostrar cómo la dinámica admite suficientes constantes de movimiento para cerrar completamente el sistema.
\end{importante}

\section{Primera cuadratura: momento angular relativo específico}
\subsection{Obtención mediante factor integrante}
Multiplicamos vectorialmente la ecuación de movimiento por el vector posición \(\vec{r}\) por la izquierda:
\[
\vec{r} \times \ddot{\vec{r}} = -\frac{\mu}{r^3}(\vec{r} \times \vec{r}) = \vec{0}
\]

Aplicamos la regla de derivación del producto vectorial:
\[
\frac{d}{dt}(\vec{r} \times \dot{\vec{r}}) = \dot{\vec{r}} \times \dot{\vec{r}} + \vec{r} \times \ddot{\vec{r}} = \vec{0}
\]

Al integrar respecto al tiempo, obtenemos una constante vectorial:
\[
\vec{h} \equiv \vec{r} \times \vec{v} = \text{cte.}
\]

Esta integral aporta \textbf{3 cuadraturas} (una por cada componente cartesiana). El vector \(\vec{h}\) es el \textbf{momento angular relativo específico}, es decir, por unidad de masa reducida.

\subsection{Interpretación física y geométrica}
\begin{itemize}
    \item \textbf{Movimiento plano:} como \(\vec{h}\) es constante y perpendicular tanto a \(\vec{r}\) como a \(\vec{v}\), el movimiento relativo queda confinado a un plano fijo perpendicular a \(\vec{h}\).
    \item \textbf{Coordenadas polares:} en el plano orbital, su módulo se expresa como \(h = r^2\dot{\theta}\).
    \item \textbf{Velocidad areal y segunda ley de Kepler:} el área infinitesimal barrida por el radio vector es \(dA = \frac{1}{2}r^2 d\theta\). Derivando respecto al tiempo:
\end{itemize}

\[
\frac{dA}{dt} = \frac{1}{2}r^2\dot{\theta} = \frac{h}{2} = \text{cte.}
\]

Esto demuestra que la velocidad areal es constante. La magnitud de \(h\) controla directamente la tasa de barrido del área, lo cual constituye la formulación dinámica de la segunda ley de Kepler.

\begin{importante}
La conservación del momento angular específico no solo reduce el número de grados efectivos del problema, sino que además garantiza que toda órbita kepleriana es plana.
\end{importante}

\section{Segunda cuadratura: energía relativa específica}
\subsection{Obtención mediante producto escalar}
Multiplicamos escalarmente la ecuación de movimiento por la velocidad \(\dot{\vec{r}}\):
\[
\dot{\vec{r}} \cdot \ddot{\vec{r}} = -\frac{\mu}{r^3}\vec{r} \cdot \dot{\vec{r}}
\]

Identificamos las derivadas exactas:
\[
\frac{1}{2}\frac{d}{dt}(v^2) = -\mu \frac{d}{dt}\left(\frac{1}{r}\right)
\]

Reordenando e integrando:
\[
\frac{d}{dt}\left(\frac{v^2}{2} - \frac{\mu}{r}\right) = 0
\]

Definimos la \textbf{energía relativa específica} como la constante resultante:
\[
\mathcal{E} \equiv \frac{v^2}{2} - \frac{\mu}{r} = \text{cte.}
\]

Esta integral escalar aporta \textbf{1 cuadratura} adicional.

\subsection{Interpretación física}
Representa la conservación de la energía mecánica por unidad de masa reducida. Su valor algebraico clasifica la trayectoria:
\begin{itemize}
    \item \(\mathcal{E} < 0\): órbita ligada (elipse).
    \item \(\mathcal{E} = 0\): caso límite de escape (parábola).
    \item \(\mathcal{E} > 0\): órbita no ligada (hipérbola).
\end{itemize}

\begin{importante}
El signo de \(\mathcal{E}\) permite distinguir si el movimiento permanece acotado o si el cuerpo escapa al infinito, incluso antes de resolver explícitamente la trayectoria.
\end{importante}

\section{Tercera cuadratura: vector de excentricidad}
\subsection{Construcción del factor integrante}
Para completar el sistema, cruzamos la ecuación de movimiento con el vector constante \(\vec{h}\):
\[
\ddot{\vec{r}} \times \vec{h} = -\frac{\mu}{r^3}\vec{r} \times (\vec{r} \times \vec{v})
\]

Aplicamos la identidad del triple producto vectorial (BAC-CAB):
\[
\vec{r} \times (\vec{r} \times \vec{v}) = \vec{r}(\vec{r} \cdot \vec{v}) - \vec{v}(\vec{r} \cdot \vec{r}) = r\dot{r}\vec{r} - r^2\vec{v}
\]

Sustituyendo y simplificando:
\[
\ddot{\vec{r}} \times \vec{h} = \mu\left(\frac{\vec{v}}{r} - \frac{\dot{r}\vec{r}}{r^2}\right)
\]

El término entre paréntesis coincide con la derivada temporal del vector unitario radial \(\hat{r} = \vec{r}/r\):
\[
\frac{d}{dt}\left(\frac{\vec{r}}{r}\right) = \frac{\dot{\vec{r}}}{r} - \frac{\vec{r}\dot{r}}{r^2}
\]

Por tanto, la ecuación se reduce a una derivada total exacta:
\[
\frac{d}{dt}(\dot{\vec{r}} \times \vec{h}) = \frac{d}{dt}\left(\mu\frac{\vec{r}}{r}\right)
\]

Integrando, obtenemos una nueva constante vectorial:
\[
\vec{e} \equiv \frac{\vec{v} \times \vec{h}}{\mu} - \frac{\vec{r}}{r} = \text{cte.}
\]

\subsection{Propiedades del vector de excentricidad}
Este vector, también llamado \textbf{vector de Laplace} o \textbf{vector de Runge--Lenz}, posee varias propiedades fundamentales:
\begin{itemize}
    \item yace en el plano orbital,
    \item es perpendicular a \(\vec{h}\),
    \item apunta hacia el periastro,
    \item su módulo coincide con la excentricidad orbital: \(e = |\vec{e}|\).
\end{itemize}

Formalmente aporta \textbf{3 cuadraturas}, aunque no todas son independientes de las integrales anteriores.

\begin{importante}
El vector de excentricidad contiene información geométrica fina sobre la orientación y la forma de la órbita. Mientras \(\vec{h}\) fija el plano orbital, \(\vec{e}\) fija la dirección del periastro dentro de ese plano.
\end{importante}

\section{Cierre del sistema: conteo de cuadraturas independientes}
Sumando las constantes obtenidas, tenemos:
\begin{itemize}
    \item momento angular relativo específico \(\vec{h}\): 3 componentes,
    \item energía relativa específica \(\mathcal{E}\): 1 escalar,
    \item vector de excentricidad \(\vec{e}\): 3 componentes.
\end{itemize}

El total aparente es de \(7\) ecuaciones. Sin embargo, el sistema no posee siete grados de libertad independientes, porque existen relaciones algebraicas que ligan estas constantes:
\begin{enumerate}
    \item ortogonalidad geométrica:
\end{enumerate}
\[
\vec{h} \cdot \vec{e} = 0
\]
\begin{enumerate}[resume]
    \item relación escalar derivada de \(\vec{e} \cdot \vec{e}\):
\end{enumerate}
\[
 e^2 = 1 + \frac{2\mathcal{E}h^2}{\mu^2}
\]

Al restar estas dependencias funcionales, el sistema queda con exactamente \textbf{6 cuadraturas independientes}, que coinciden con el número de variables dinámicas iniciales. Esto demuestra que el problema relativo de dos cuerpos es \textbf{completamente integrable} y permite derivar la ecuación de la trayectoria \(r(\theta)\) sin recurrir, en principio, a métodos numéricos.

\section{Síntesis}
En esta clase se obtuvo el conjunto fundamental de integrales primeras del problema relativo de dos cuerpos:
\begin{itemize}
    \item conservación del momento angular específico,
    \item conservación de la energía específica,
    \item conservación del vector de excentricidad.
\end{itemize}

Estas cuadraturas no solo permiten cerrar analíticamente el sistema, sino que preparan el camino para la deducción explícita de la ecuación de la cónica orbital y el estudio detallado de los elementos orbitales en la clase siguiente.
\chapter{Clase 17: Solución analítica del problema relativo y plano de Laplace}

Esta clase continúa el cierre analítico del problema relativo de dos cuerpos. A partir de las cuadraturas obtenidas en la clase anterior, se reduce el problema al plano orbital y se deriva la ecuación polar de la trayectoria. El objetivo central es mostrar que la solución relativa corresponde siempre a una sección cónica y que su geometría queda completamente determinada por las constantes del movimiento.

\section{Reducción al plano de la trayectoria}
Partimos del conteo de cuadraturas obtenido en el espacio tridimensional:
\begin{itemize}
    \item momento angular relativo específico \(\vec{h}\): 3 componentes,
    \item energía relativa específica \(\mathcal{E}\): 1 escalar,
    \item vector de Laplace o de excentricidad \(\vec{e}\): 3 componentes.
\end{itemize}

En apariencia esto produce 7 constantes, pero una de ellas queda ligada por la relación
\[
\vec{h} \cdot \vec{e} = 0,
\]
por lo que el sistema tiene finalmente \textbf{6 cuadraturas independientes}.

Al pasar al \textbf{plano de la trayectoria} —también llamado plano orbital o plano perpendicular a \(\vec{h}\)— el problema se reduce a dos dimensiones espaciales. En ese plano, el número de variables dinámicas independientes es \(4\): \((x, y, \dot{x}, \dot{y})\). En consecuencia, basta con describir el movimiento mediante cuatro constantes escalares equivalentes:
\begin{enumerate}
    \item \(h = |\vec{h}|\), módulo del momento angular específico,
    \item \(\mathcal{E}\), energía relativa específica,
    \item la dirección de \(\vec{e}\) en el plano orbital,
    \item el módulo \(e = |\vec{e}|\).
\end{enumerate}

Este marco permite enfocarse exclusivamente en la geometría y dinámica plana de la órbita.

\begin{importante}
La reducción al plano orbital no es una aproximación. Es una consecuencia exacta de la conservación del momento angular: toda la dinámica relativa ocurre en un único plano fijo.
\end{importante}

\section{Análisis del vector de Laplace}
El vector de excentricidad se define como:
\[
\vec{e} = \frac{\vec{v} \times \vec{h}}{\mu} - \frac{\vec{r}}{r}
\]

\subsection{Dirección y significado geométrico}
Sus propiedades más importantes son:
\begin{itemize}
    \item \(\vec{e}\) es constante en magnitud y dirección,
    \item está contenido en el plano orbital, pues es combinación lineal de \(\vec{r}\) y \(\vec{v}\),
    \item apunta hacia el \textbf{periastro}, es decir, el punto de mínima distancia radial.
\end{itemize}

El ángulo \(\theta\) medido desde la dirección de \(\vec{e}\) hasta \(\vec{r}\) se denomina \textbf{verdadera anomalía}. Esta elección angular es especialmente natural, porque toma como referencia el eje geométricamente privilegiado de la órbita.

\subsection{Cálculo de la magnitud \texorpdfstring{$e^2$}{e²}}
Calculamos el producto escalar \(\vec{e} \cdot \vec{e}\) para obtener una expresión analítica en función de las integrales del movimiento:
\[
 e^2 = \left( \frac{\vec{v} \times \vec{h}}{\mu} - \frac{\vec{r}}{r} \right) \cdot \left( \frac{\vec{v} \times \vec{h}}{\mu} - \frac{\vec{r}}{r} \right)
\]

Desarrollando:
\[
 e^2 = \frac{|\vec{v} \times \vec{h}|^2}{\mu^2} - \frac{2}{\mu r}(\vec{v} \times \vec{h}) \cdot \vec{r} + 1
\]

\textbf{Paso A. Primer término: identidad de Lagrange}
\[
 |\vec{v} \times \vec{h}|^2 = v^2 h^2 - (\vec{v} \cdot \vec{h})^2
\]
Como \(\vec{h} = \vec{r} \times \vec{v}\), entonces \(\vec{v} \cdot \vec{h} = 0\). Por tanto:
\[
 |\vec{v} \times \vec{h}|^2 = v^2 h^2
\]

\textbf{Paso B. Segundo término: propiedad cíclica del triple producto escalar}
\[
 (\vec{v} \times \vec{h}) \cdot \vec{r} = \vec{h} \cdot (\vec{r} \times \vec{v}) = \vec{h} \cdot \vec{h} = h^2
\]

\textbf{Paso C. Ensamblaje y sustitución de la energía}
\[
 e^2 = \frac{v^2 h^2}{\mu^2} - \frac{2 h^2}{\mu r} + 1
\]
\[
 e^2 = 1 + \frac{2h^2}{\mu^2}\left(\frac{v^2}{2} - \frac{\mu}{r}\right)
\]

Reconociendo que
\[
\mathcal{E} = \frac{v^2}{2} - \frac{\mu}{r},
\]
obtenemos la relación fundamental:
\[
\boxed{e^2 = 1 + \frac{2\mathcal{E}h^2}{\mu^2}}
\]

\begin{importante}
Esta ecuación conecta directamente la energía del sistema con la forma geométrica de la órbita. La excentricidad no es un parámetro puramente geométrico aislado: está determinada por las constantes dinámicas del movimiento.
\end{importante}

\subsection{Clasificación de órbitas según la energía}
La relación anterior permite clasificar las cónicas orbitales:
\begin{itemize}
    \item si \(\mathcal{E} < 0\), entonces \(0 \le e < 1\): la órbita es una \textbf{elipse},
    \item si \(\mathcal{E} = 0\), entonces \(e = 1\): la órbita es una \textbf{parábola},
    \item si \(\mathcal{E} > 0\), entonces \(e > 1\): la órbita es una \textbf{hipérbola}.
\end{itemize}

Así, la energía específica determina si la trayectoria es ligada, marginal o de escape.

\section{Deducción de la ecuación polar de la trayectoria}
Proyectamos ahora el vector de Laplace sobre el vector posición \(\vec{r}\) para obtener la relación funcional entre \(r\) y \(\theta\).

\subsection{Proyección geométrica}
Por definición del producto punto:
\[
\vec{e} \cdot \vec{r} = e r \cos\theta
\]
donde \(\theta\) es la verdadera anomalía, es decir, el ángulo entre \(\vec{e}\) y \(\vec{r}\).

\subsection{Proyección algebraica}
Usando la definición del vector \(\vec{e}\):
\[
\vec{e} \cdot \vec{r} = \left(\frac{\vec{v} \times \vec{h}}{\mu} - \frac{\vec{r}}{r}\right) \cdot \vec{r}
\]
\[
\vec{e} \cdot \vec{r} = \frac{(\vec{v} \times \vec{h}) \cdot \vec{r}}{\mu} - \frac{\vec{r} \cdot \vec{r}}{r}
\]

Usando que \((\vec{v} \times \vec{h}) \cdot \vec{r} = h^2\) y que \(\vec{r} \cdot \vec{r} = r^2\), resulta:
\[
\vec{e} \cdot \vec{r} = \frac{h^2}{\mu} - r
\]

\subsection{Igualación y despeje de \texorpdfstring{$r(\theta)$}{r(theta)}}
Igualando la expresión geométrica con la algebraica:
\[
 e r \cos\theta = \frac{h^2}{\mu} - r
\]
\[
 r(1 + e\cos\theta) = \frac{h^2}{\mu}
\]

Definiendo el \textbf{parámetro} o \textbf{semilatus rectum}
\[
 p \equiv \frac{h^2}{\mu},
\]
se llega a la ecuación polar final de la trayectoria:
\[
\boxed{r(\theta) = \frac{p}{1 + e\cos\theta}}
\]

Esta es la ecuación general de una sección cónica con un foco en el origen del sistema relativo.

\begin{importante}
La ecuación \(r(\theta) = \frac{p}{1 + e\cos\theta}\) muestra que el problema relativo de dos cuerpos tiene como soluciones exactas cónicas: circunferencias, elipses, parábolas e hipérbolas aparecen como casos particulares según el valor de \(e\).
\end{importante}

\section{Resumen de la integración completa}
Los resultados principales de esta clase pueden sintetizarse así:
\begin{enumerate}
    \item El problema relativo de dos cuerpos posee suficientes cuadraturas independientes para ser completamente integrable.
    \item En el plano orbital, la dinámica queda caracterizada por \(h\), \(\mathcal{E}\), la dirección de \(\vec{e}\) y su magnitud \(e\).
    \item El vector \(\vec{e}\) fija la orientación de la órbita y su módulo determina la excentricidad geométrica mediante
\[
 e^2 = 1 + \frac{2\mathcal{E}h^2}{\mu^2}.
\]
    \item La proyección de \(\vec{e}\) sobre \(\vec{r}\) conduce directamente a
\[
 r(\theta) = \frac{p}{1 + e\cos\theta},
\]
lo que demuestra analíticamente que la trayectoria relativa es siempre una cónica.
\end{enumerate}

Este desarrollo cierra la solución analítica del problema de dos cuerpos y deja listo el marco conceptual para estudiar con más detalle la geometría orbital, los elementos de la cónica y su representación en distintos sistemas de referencia.
\chapter{Clase 18: Geometría de las cónicas (parte 1)}

Esta clase inicia el estudio geométrico de las cónicas desde la perspectiva de la mecánica celeste. Después de haber deducido la ecuación polar de la trayectoria relativa, el siguiente paso es comprender en qué sistema de referencia esa ecuación adopta su forma más simple y cómo se transforma al sistema de referencia del observador. El objetivo es establecer con claridad el \textbf{sistema natural de la cónica} y desarrollar la herramienta algebraica básica para relacionarlo con un sistema externo: la \textbf{rotación plana}.

\section{Introducción y repaso del marco teórico}
El problema relativo de dos cuerpos, una vez integradas sus cuadraturas, conduce a la ecuación de la trayectoria en coordenadas polares:
\[
 r(\theta) = \frac{p}{1 + e\cos\theta}
\]
donde \(p = h^2/\mu\) es el parámetro semilatus rectum, \(e = |\vec{e}|\) es la excentricidad y \(\theta\) es la verdadera anomalía medida desde la dirección del vector de Laplace \(\vec{e}\).

Esta expresión describe completamente la forma de la órbita, pero está escrita en un sistema de coordenadas muy particular: aquel cuyo eje polar coincide con la dirección de \(\vec{e}\) y cuyo origen se sitúa en el foco gravitacional. A este marco se le denomina \textbf{sistema natural de la cónica}.

El objetivo analítico de esta clase es formalizar la geometría de las secciones cónicas en su sistema natural, identificar las limitaciones que presenta para la descripción observacional y construir rigurosamente las herramientas de rotación plana necesarias para pasar del sistema natural al sistema de referencia del observador.

\begin{importante}
La ecuación polar de la cónica no depende solo de la forma geométrica de la órbita, sino también de la elección del sistema de referencia. Su forma simple se obtiene únicamente en el sistema natural alineado con la órbita.
\end{importante}

\section{El sistema natural de la cónica}
Toda sección cónica posee una geometría intrínseca que se describe de manera más simple cuando se adopta un sistema de ejes alineado con sus elementos de simetría.

\subsection{Definición del marco natural \texorpdfstring{$(x',y')$}{(x',y')}}
Se define un sistema cartesiano plano \((x', y')\) con las siguientes características:
\begin{itemize}
    \item el eje \(x'\) coincide con el \textbf{eje de simetría} de la cónica,
    \item el origen puede ubicarse en el \textbf{foco} (sistema focal) o en el \textbf{centro geométrico} (sistema central), dependiendo de la conveniencia analítica,
    \item el eje \(y'\) es perpendicular a \(x'\) y completa la base ortonormal derecha del plano orbital.
\end{itemize}

En este sistema, la descripción analítica se simplifica notablemente porque la cónica presenta simetría respecto a \(x'\) y \(y'\).

\subsection{Ecuaciones en el sistema natural}
Dos descripciones típicas en este sistema son:
\begin{itemize}
    \item \textbf{Ecuación polar con origen en el foco:}
\[
 r = \frac{p}{1 + e\cos\theta'}
\]
    donde \(\theta'\) es el ángulo medido desde el semieje positivo \(x'\) hasta el vector posición \(\vec{r}\).
    \item \textbf{Ecuaciones paramétricas cartesianas:} dependiendo del tipo de cónica, las coordenadas \((x', y')\) pueden expresarse en función de un parámetro auxiliar. Por ejemplo, para la elipse usando la anomalía excéntrica \(E\):
\[
 x' = a(\cos E - e), \quad y' = a\sqrt{1-e^2}\sin E
\]
\end{itemize}

Estas expresiones conservan su forma funcional dentro del sistema natural, independientemente de cómo esté orientada la órbita en el espacio físico.

\section{El sistema de referencia del observador y el problema de alineación}
En la práctica astronómica, las posiciones y velocidades se miden respecto a un sistema de referencia fijo o inercial, denominado \textbf{sistema del observador} \((x,y,z)\). Este sistema puede ser, por ejemplo:
\begin{itemize}
    \item el plano ecuatorial terrestre,
    \item el plano de la eclíptica,
    \item el plano tangente a la esfera celeste en la dirección de un objeto.
\end{itemize}

El problema fundamental es que casi nunca el sistema de ejes utilizado en las observaciones coincide con el sistema natural de la órbita. El plano orbital puede estar inclinado respecto al plano de referencia y la línea de ápsides, es decir, la dirección de \(\vec{e}\), puede formar un ángulo arbitrario con el eje \(x\) del observador.

Por lo tanto, para comparar predicciones teóricas con datos observacionales, es necesario establecer una transformación matemática rigurosa que relacione las coordenadas del sistema natural \((x',y')\) con las del sistema del observador \((x,y)\). Como el movimiento sigue siendo plano en esta etapa, la transformación se reduce inicialmente a una \textbf{rotación en el plano}.

\begin{importante}
La dificultad observacional no está en la ecuación de la órbita, sino en que los datos rara vez están expresados en el sistema donde esa ecuación adopta su forma más simple.
\end{importante}

\section{Rotación de la cónica en el plano}
Consideremos dos sistemas cartesianos ortonormales que comparten el mismo origen:
\begin{itemize}
    \item sistema del observador: base \(\{\hat{i}, \hat{j}\}\),
    \item sistema natural de la cónica: base \(\{\hat{i}', \hat{j}'\}\).
\end{itemize}

El sistema natural está rotado un ángulo \(\theta\) respecto al sistema del observador, en sentido antihorario. El vector posición \(\vec{r}\) es el mismo objeto geométrico, pero sus componentes cambian con la base elegida:
\[
\vec{r} = x\hat{i} + y\hat{j} = x'\hat{i}' + y'\hat{j}'
\]

\subsection{Relación entre vectores unitarios}
Para expresar los componentes en una base en función de los de la otra, proyectamos los vectores unitarios rotados sobre la base fija:
\[
\hat{i}' = \cos\theta\,\hat{i} + \sin\theta\,\hat{j}
\]
\[
\hat{j}' = -\sin\theta\,\hat{i} + \cos\theta\,\hat{j}
\]

Sustituyendo en \(\vec{r} = x'\hat{i}' + y'\hat{j}'\) y agrupando términos respecto a \(\hat{i}\) y \(\hat{j}\), obtenemos:
\[
 x\hat{i} + y\hat{j} = (x'\cos\theta - y'\sin\theta)\hat{i} + (x'\sin\theta + y'\cos\theta)\hat{j}
\]

Igualando componentes, se deduce la \textbf{regla de rotación directa}:
\[
\begin{cases}
 x = x'\cos\theta - y'\sin\theta \\
 y = x'\sin\theta + y'\cos\theta
\end{cases}
\]

\section{Derivación analítica de la matriz de rotación \texorpdfstring{$R_z(\theta)$}{Rz(theta)}}
El sistema anterior puede escribirse de manera compacta mediante álgebra matricial. Definimos los vectores columna de coordenadas:
\[
\begin{pmatrix} x \\ y \end{pmatrix} = \mathbf{R}_z(\theta)\begin{pmatrix} x' \\ y' \end{pmatrix}
\]

Aquí, \(\mathbf{R}_z(\theta)\) es la \textbf{matriz de rotación en el plano \(x\)-\(y\) alrededor del eje \(z\)}:
\[
\mathbf{R}_z(\theta) =
\begin{pmatrix}
\cos\theta & -\sin\theta \\
\sin\theta & \cos\theta
\end{pmatrix}
\]

\subsection{Propiedades analíticas de \texorpdfstring{$\mathbf{R}_z(\theta)$}{Rz(theta)}}
Esta matriz posee propiedades fundamentales:
\begin{enumerate}
    \item \textbf{Ortogonalidad:}
\[
\mathbf{R}_z^T(\theta)\mathbf{R}_z(\theta) = \mathbf{I}
\]
    lo cual garantiza que la transformación preserva productos escalares, distancias y ángulos.
    \item \textbf{Determinante unitario:}
\[
\det(\mathbf{R}_z) = \cos^2\theta + \sin^2\theta = 1
\]
    indicando que la rotación es propia y no incluye reflexiones.
    \item \textbf{Preservación de la norma:}
\[
 x^2 + y^2 = x'^2 + y'^2
\]
    consistente con la invariancia de la distancia radial \(r = |\vec{r}|\).
\end{enumerate}

Esta matriz es la herramienta fundamental para trasladar cualquier descripción geométrica de la cónica desde su sistema natural al sistema del observador.

\section{La rotación inversa y el cambio de signo angular}
Hasta ahora se ha trabajado en el caso plano, donde bastaba una sola rotación alrededor del eje \(z\). Sin embargo, para describir una órbita real en un sistema inercial del observador es necesario introducir explícitamente las tres coordenadas \((x,y,z)\) y la orientación espacial completa del plano orbital.

La idea central sigue siendo la misma: el sistema natural de la órbita es el sistema donde la cónica tiene la forma más simple. Lo que cambia es que ahora ese sistema debe incrustarse dentro de un espacio tridimensional.

\subsection{Rotación inversa en el caso plano como antesala del caso espacial}
En dos dimensiones, la transformación inversa se obtenía mediante
\[
\begin{pmatrix} x' \\ y' \end{pmatrix} = \mathbf{R}_z^{-1}(\theta)\begin{pmatrix} x \\ y \end{pmatrix}
\]
y, como \(\mathbf{R}_z\) es ortogonal,
\[
\mathbf{R}_z^{-1}(\theta) = \mathbf{R}_z^T(\theta) = \mathbf{R}_z(-\theta).
\]

Este resultado sigue siendo válido en tres dimensiones para cualquier matriz de rotación elemental: la transformación inversa se obtiene invirtiendo el orden de las rotaciones y cambiando el signo de los ángulos.

\subsection{Unicidad de la transformación en el espacio}
En el problema tridimensional no aparecen matrices distintas para posición y para velocidad. La misma transformación geométrica que orienta el sistema perifocal dentro del espacio inercial actúa sobre cualquier vector del plano orbital. Por ello, una vez construida la matriz total de rotación, esta sirve tanto para transformar \(\vec{r}\) como \(\vec{v}\), o cualquier otro vector expresado en el sistema natural.

\section{Síntesis y preparación para la tridimensionalidad}
Los resultados obtenidos hasta aquí muestran que la descripción plana de la cónica es solamente una etapa intermedia. Para conectar la órbita con el sistema del observador es necesario reconstruir la geometría en tres dimensiones usando coordenadas \((x,y,z)\).

\subsection{7.4 Deducción tridimensional en coordenadas \texorpdfstring{$x,y,z$}{x,y,z}}
Introducimos el sistema natural tridimensional de la órbita, también llamado \textbf{sistema perifocal}, con base ortonormal \(\{\hat{p},\hat{q},\hat{w}\}\):
\begin{itemize}
    \item \(\hat{p}\) apunta hacia el periastro,
    \item \(\hat{q}\) yace en el plano orbital y es perpendicular a \(\hat{p}\),
    \item \(\hat{w}\) es paralelo al momento angular específico \(\vec{h}\).
\end{itemize}

En este sistema, la órbita sigue contenida en un plano, por lo que la coordenada perpendicular vale cero. Así, la posición relativa y la velocidad relativa se escriben como
\[
\vec{r}_{\mathrm{pf}} =
\begin{pmatrix}
 r\cos f \\
 r\sin f \\
 0
\end{pmatrix}
\]
\[
\vec{v}_{\mathrm{pf}} =
\sqrt{\frac{\mu}{p}}
\begin{pmatrix}
 -\sin f \\
 e+\cos f \\
 0
\end{pmatrix}
\]
donde \(f\) es la verdadera anomalía.

El problema consiste en expresar estos vectores en el sistema inercial del observador \(\{\hat{i},\hat{j},\hat{k}\}\), es decir, en términos de las coordenadas \((x,y,z)\). Para ello deben introducirse tres ángulos orbitales:
\begin{enumerate}
    \item la longitud del nodo ascendente \(\Omega\),
    \item la inclinación orbital \(i\),
    \item el argumento del periastro \(\omega\).
\end{enumerate}

La construcción geométrica se hace en tres pasos.

\textbf{Paso 1. Rotación por \(\omega\) alrededor de \(z\).}  
Esta rotación orienta el eje \(\hat{p}\) dentro del plano orbital. La matriz asociada es
\[
\mathbf{R}_z(\omega)=
\begin{pmatrix}
\cos\omega & -\sin\omega & 0 \\
\sin\omega & \cos\omega & 0 \\
0 & 0 & 1
\end{pmatrix}.
\]

\textbf{Paso 2. Rotación por \(i\) alrededor de \(x\).}  
Esta rotación inclina el plano orbital respecto al plano de referencia:
\[
\mathbf{R}_x(i)=
\begin{pmatrix}
1 & 0 & 0 \\
0 & \cos i & -\sin i \\
0 & \sin i & \cos i
\end{pmatrix}.
\]

\textbf{Paso 3. Rotación por \(\Omega\) alrededor de \(z\).}  
Esta última rotación fija la dirección de la línea de nodos en el sistema inercial:
\[
\mathbf{R}_z(\Omega)=
\begin{pmatrix}
\cos\Omega & -\sin\Omega & 0 \\
\sin\Omega & \cos\Omega & 0 \\
0 & 0 & 1
\end{pmatrix}.
\]

La matriz total que transforma del sistema perifocal al sistema inercial es entonces
\[
\mathbf{Q} = \mathbf{R}_z(\Omega)\,\mathbf{R}_x(i)\,\mathbf{R}_z(\omega).
\]
Por tanto,
\[
\vec{r}_{\mathrm{I}} = \mathbf{Q}\,\vec{r}_{\mathrm{pf}},
\qquad
\vec{v}_{\mathrm{I}} = \mathbf{Q}\,\vec{v}_{\mathrm{pf}}.
\]

Al desarrollar el producto matricial y aplicar la matriz a la posición perifocal, se obtiene la deducción explícita en coordenadas cartesianas tridimensionales:
\[
\begin{pmatrix}
 x \\
 y \\
 z
\end{pmatrix}
=
\mathbf{Q}
\begin{pmatrix}
 r\cos f \\
 r\sin f \\
 0
\end{pmatrix}
=
 r
\begin{pmatrix}
 \cos\Omega\cos(\omega+f)-\sin\Omega\sin(\omega+f)\cos i \\
 \sin\Omega\cos(\omega+f)+\cos\Omega\sin(\omega+f)\cos i \\
 \sin(\omega+f)\sin i
\end{pmatrix}.
\]

De aquí se leen directamente las coordenadas espaciales de la órbita:
\[
 x = r\bigl[\cos\Omega\cos(\omega+f)-\sin\Omega\sin(\omega+f)\cos i\bigr]
\]
\[
 y = r\bigl[\sin\Omega\cos(\omega+f)+\cos\Omega\sin(\omega+f)\cos i\bigr]
\]
\[
 z = r\sin(\omega+f)\sin i.
\]

Estas expresiones constituyen la forma explícita de la órbita kepleriana en el espacio tridimensional del observador.

\begin{importante}
La coordenada \(z\) aparece únicamente porque el plano orbital está inclinado. Si \(i=0\), entonces \(z=0\) y se recupera exactamente el caso plano estudiado antes.
\end{importante}

\subsection*{7.5 Velocidad tridimensional en el sistema inercial}
La misma matriz \(\mathbf{Q}\) transforma la velocidad desde el sistema perifocal al sistema inercial. Por tanto,
\[
\begin{pmatrix}
 \dot{x} \\
 \dot{y} \\
 \dot{z}
\end{pmatrix}
=
\mathbf{Q}
\sqrt{\frac{\mu}{p}}
\begin{pmatrix}
 -\sin f \\
 e+\cos f \\
 0
\end{pmatrix}.
\]

Desarrollando componente a componente:
\[
\dot{x} = \sqrt{\frac{\mu}{p}}\left[-\cos\Omega\sin(\omega+f)-\sin\Omega\cos(\omega+f)\cos i + e\bigl(-\cos\Omega\sin\omega-\sin\Omega\cos\omega\cos i\bigr)\right]
\]
\[
\dot{y} = \sqrt{\frac{\mu}{p}}\left[-\sin\Omega\sin(\omega+f)+\cos\Omega\cos(\omega+f)\cos i + e\bigl(-\sin\Omega\sin\omega+\cos\Omega\cos\omega\cos i\bigr)\right]
\]
\[
\dot{z} = \sqrt{\frac{\mu}{p}}\left[\cos(\omega+f)\sin i + e\cos\omega\sin i\right].
\]

Así, la órbita queda completamente descrita en el sistema del observador mediante las seis componentes \((x,y,z,\dot{x},\dot{y},\dot{z})\).

\subsection*{7.6 Interpretación geométrica final}
La dinámica kepleriana sigue siendo plana en su naturaleza: el cuerpo se mueve siempre en un plano fijo perpendicular a \(\vec{h}\). Sin embargo, para un observador inercial, ese plano puede encontrarse arbitrariamente orientado en el espacio. La tridimensionalidad no introduce una nueva fuerza ni una nueva ecuación orbital; introduce una nueva \textbf{geometría de orientación}.

En consecuencia:
\begin{itemize}
    \item la ecuación polar fija la forma de la órbita en su plano natural,
    \item los ángulos \(\Omega\), \(i\) y \(\omega\) fijan la orientación del plano y del periastro en el espacio,
    \item las fórmulas de \((x,y,z)\) y \((\dot{x},\dot{y},\dot{z})\) permiten comparar directamente la teoría con observaciones y simulaciones numéricas.
\end{itemize}

Este es el puente completo entre la solución analítica del problema de dos cuerpos y su representación en un sistema inercial tridimensional.

\begin{importante}
La geometría de las cónicas no solo determina la forma de la trayectoria orbital. También fija la estructura algebraica de las transformaciones de coordenadas que hacen posible comparar teoría, simulación y observación.
\end{importante}
\chapter{Clase 20: Elementos orbitales en el problema de dos cuerpos}

Esta clase cierra el puente entre la solución analítica del problema de dos cuerpos y su uso práctico en astronomía y astrodinámica. Después de haber formulado la órbita en el sistema natural y de haber introducido las rotaciones necesarias para orientarla en el espacio, el siguiente paso es condensar toda esa información en un conjunto mínimo de parámetros geométricos y dinámicos: los \textbf{elementos orbitales}.

\section{Repaso: sistema natural y rotación plana}
La ecuación de la trayectoria relativa en coordenadas polares se expresa de forma canónica en el \textbf{sistema natural de la cónica}, también llamado sistema perifocal, cuyo eje \(x'\) coincide con la dirección del vector de Laplace \(\vec{e}\) y cuyo origen se sitúa en el foco gravitacional:
\[
 r = \frac{p}{1 + e\cos f}
\]
donde \(p = h^2/\mu\) es el semilatus rectum, \(e = |\vec{e}|\) la excentricidad y \(f\) la verdadera anomalía.

En este mismo sistema, las coordenadas cartesianas se obtienen proyectando \(\vec{r}\) sobre la base rotada \(\{\hat{i}', \hat{j}'\}\):
\[
 x' = r\cos f, \quad y' = r\sin f
\]

El \textbf{sistema del observador} \(\{\hat{i}, \hat{j}, \hat{k}\}\) rara vez coincide con el sistema natural. En el caso plano, la orientación relativa se describe mediante una rotación dada por
\[
\begin{pmatrix} x \\ y \end{pmatrix} =
\begin{pmatrix}
\cos\omega & -\sin\omega \\
\sin\omega & \cos\omega
\end{pmatrix}
\begin{pmatrix} x' \\ y' \end{pmatrix}
\]
Esta idea es la base para extender el análisis al espacio tridimensional.

\begin{importante}
Los elementos orbitales surgen precisamente como una manera compacta de describir la forma, la orientación y la posición instantánea del cuerpo dentro de su órbita kepleriana.
\end{importante}

\section{La pesadilla de los cuadrantes}
En la determinación analítica de ángulos orbitales, el uso directo de funciones trigonométricas inversas como \(\arccos\) introduce una ambigüedad fundamental: su rango principal está restringido a \([0,\pi]\). Esto impide distinguir entre configuraciones geométricas simétricas respecto al eje de referencia.

\subsection{Ambigüedad en la longitud del nodo ascendente \texorpdfstring{$\Omega$}{Omega}}
El vector nodo \(\vec{n}\) se define como la intersección del plano orbital con el plano de referencia, proyectada sobre dicho plano:
\[
\vec{n} = \hat{k} \times \vec{h}
\]

Si calculamos
\[
\Omega = \arccos\left(\frac{n_x}{n}\right),
\]
obtenemos un valor en \([0,\pi]\). Sin embargo, si \(n_y < 0\), el nodo ascendente se encuentra en el semiplano inferior, por lo que el ángulo real debe corregirse:
\[
\Omega =
\begin{cases}
\arccos(n_x/n), & \text{si } n_y \ge 0 \\
2\pi - \arccos(n_x/n), & \text{si } n_y < 0
\end{cases}
\]

\subsection{Ambigüedad en el argumento del periapsis \texorpdfstring{$\omega$}{omega}}
De manera análoga,
\[
\omega = \arccos\left(\frac{\vec{n}\cdot\vec{e}}{ne}\right)
\]
devuelve un valor principal. La corrección de cuadrante depende del signo de la componente \(z\) del vector de excentricidad, o equivalentemente del sentido relativo entre \(\vec{n}\), \(\vec{e}\) y \(\vec{h}\):
\[
\omega =
\begin{cases}
\arccos\left(\frac{\vec{n}\cdot\vec{e}}{ne}\right), & \text{si } e_z \ge 0 \\
2\pi - \arccos\left(\frac{\vec{n}\cdot\vec{e}}{ne}\right), & \text{si } e_z < 0
\end{cases}
\]

\begin{importante}
La determinación unívoca de ángulos orbitales requiere usar \(\operatorname{atan2}(y,x)\) o, de forma equivalente, corregir los cuadrantes inspeccionando el signo de las componentes transversales adecuadas.
\end{importante}

\section{Extensión a tres dimensiones: la cónica levantada}
Cuando el plano orbital no coincide con el plano de referencia, hacen falta tres rotaciones sucesivas para alinear el sistema perifocal con el sistema del observador. Esta secuencia corresponde a los \textbf{ángulos de Euler} en la convención \(3\)-\(1\)-\(3\), ampliamente usada en mecánica celeste.

\subsection{Línea de nodos y referencia planetaria}
Introducimos tres objetos geométricos clave:
\begin{itemize}
    \item el \textbf{plano de referencia}, usualmente el plano ecuatorial terrestre o el plano de la eclíptica,
    \item la \textbf{línea de nodos}, intersección entre el plano orbital y el plano de referencia,
    \item el \textbf{nodo ascendente}, punto donde el cuerpo cruza el plano de referencia moviéndose hacia el hemisferio norte \((z>0)\).
\end{itemize}

\subsection{Definición de los ángulos de orientación}
Los ángulos espaciales fundamentales son:
\begin{enumerate}
    \item \textbf{Longitud del nodo ascendente \(\Omega\):} ángulo en el plano de referencia desde el eje \(x\) de referencia hasta la línea de nodos.
    \item \textbf{Inclinación orbital \(I\):} ángulo entre el vector momento angular \(\vec{h}\) y el eje \(z\) de referencia.
    \item \textbf{Argumento del periapsis \(\omega\):} ángulo medido dentro del plano orbital desde la línea de nodos hasta la dirección del periapsis \((\vec{e})\).
\end{enumerate}

Estos tres ángulos, junto con \(p\), \(e\) y \(f\), determinan completamente la orientación y forma de la órbita en el espacio tridimensional.

\section{Matriz de rotación de Euler en 3D}
La transformación desde el sistema perifocal al sistema del observador se realiza mediante el producto de tres matrices de rotación elementales. La convención estándar aplica las rotaciones en el orden:
\begin{enumerate}
    \item \(\mathbf{R}_z(\omega)\): alinea el eje perifocal \(x'\) con la línea de nodos,
    \item \(\mathbf{R}_x(I)\): inclina el plano orbital respecto al plano de referencia,
    \item \(\mathbf{R}_z(\Omega)\): alinea la línea de nodos con el eje \(x\) del observador.
\end{enumerate}

La matriz de rotación de Euler resultante es:
\[
\mathbf{R}_{\text{Euler}} = \mathbf{R}_z(\Omega)\,\mathbf{R}_x(I)\,\mathbf{R}_z(\omega)
\]
Desarrollando el producto matricial:
\[
\mathbf{R}_{\text{Euler}} =
\begin{pmatrix}
\cos\Omega\cos\omega - \sin\Omega\sin\omega\cos I & -\cos\Omega\sin\omega - \sin\Omega\cos\omega\cos I & \sin\Omega\sin I \\
\sin\Omega\cos\omega + \cos\Omega\sin\omega\cos I & -\sin\Omega\sin\omega + \cos\Omega\cos\omega\cos I & -\cos\Omega\sin I \\
\sin\omega\sin I & \cos\omega\sin I & \cos I
\end{pmatrix}
\]

La posición en el sistema del observador se obtiene como
\[
\vec{r}_{\text{obs}} = \mathbf{R}_{\text{Euler}}\,\vec{r}_{\text{per}}
\]
donde
\[
\vec{r}_{\text{per}} = (r\cos f,\; r\sin f,\; 0)^T.
\]

La transformación inversa usa la transpuesta, válida por la ortogonalidad de \(\mathbf{R}_{\text{Euler}}\):
\[
\vec{r}_{\text{per}} = \mathbf{R}_{\text{Euler}}^T\,\vec{r}_{\text{obs}}
\]

\begin{importante}
La matriz de Euler resume de manera compacta toda la orientación espacial de la órbita. La forma de la cónica se mantiene intacta; lo que cambia es su orientación con respecto al observador.
\end{importante}

\section{Elementos orbitales clásicos y sus variantes}
Los seis parámetros que describen completamente una órbita kepleriana se agrupan en los \textbf{elementos orbitales clásicos}:
\begin{center}
\begin{tabular}{lll}
\textbf{Elemento} & \textbf{Símbolo} & \textbf{Significado geométrico} \\
Semilatus rectum & \(p\) & tamaño o escala de la cónica \\
Excentricidad & \(e\) & forma de la cónica \\
Inclinación & \(I\) & orientación del plano orbital \\
Longitud del nodo ascendente & \(\Omega\) & orientación de la línea de nodos \\
Argumento del periapsis & \(\omega\) & orientación del periapsis en el plano \\
Verdadera anomalía & \(f\) & posición instantánea del cuerpo
\end{tabular}
\end{center}

\subsection{Variantes descriptivas}
Dependiendo del tipo de problema, el parámetro \(p\) puede reemplazarse por otras cantidades de tamaño:
\begin{itemize}
    \item \textbf{Elementos asteroidales:} \((a, e, I, \Omega, \omega, f)\), donde \(a\) es el semieje mayor. Son útiles para órbitas cerradas \((e<1)\), con
\[
 p = a(1-e^2).
\]
    \item \textbf{Elementos cometarios:} \((q, e, I, \Omega, \omega, f)\), donde \(q\) es la distancia al periapsis. Son ventajosos para órbitas casi parabólicas \((e\approx 1)\), con
\[
 p = q(1+e).
\]
\end{itemize}

\section{Determinación de órbita: del vector de estado a los elementos}
El vector de estado \(\{\vec{r},\vec{v}\}\) contiene toda la información necesaria para recuperar los elementos orbitales. El procedimiento analítico estándar es el siguiente.

\subsection{Paso 1. Vectores fundamentales}
\[
\vec{h} = \vec{r} \times \vec{v}, \qquad
\vec{e} = \frac{1}{\mu}(\vec{v} \times \vec{h}) - \frac{\vec{r}}{r}, \qquad
\vec{n} = \hat{k} \times \vec{h}
\]

\subsection{Paso 2. Parámetros de tamaño y forma}
\[
 h = |\vec{h}|, \qquad e = |\vec{e}|, \qquad p = \frac{h^2}{\mu}
\]

\subsection{Paso 3. Orientación del plano orbital}
\[
 I = \arccos\left(\frac{h_z}{h}\right), \qquad n = |\vec{n}|
\]
\[
\Omega =
\begin{cases}
\arccos(n_x/n), & n_y \ge 0 \\
2\pi - \arccos(n_x/n), & n_y < 0
\end{cases}
\]

\subsection{Paso 4. Orientación del periapsis y posición instantánea}
\[
\omega =
\begin{cases}
\arccos\left(\frac{\vec{n}\cdot\vec{e}}{ne}\right), & e_z \ge 0 \\
2\pi - \arccos\left(\frac{\vec{n}\cdot\vec{e}}{ne}\right), & e_z < 0
\end{cases}
\]
\[
 f =
\begin{cases}
\arccos\left(\frac{\vec{e}\cdot\vec{r}}{er}\right), & \vec{r}\cdot\vec{v} \ge 0 \\
2\pi - \arccos\left(\frac{\vec{e}\cdot\vec{r}}{er}\right), & \vec{r}\cdot\vec{v} < 0
\end{cases}
\]

Este algoritmo permite pasar de \((x,y,z,v_x,v_y,v_z)\) a \((p,e,I,\Omega,\omega,f)\) de manera biunívoca, siempre que se apliquen correctamente las correcciones de cuadrante.

\begin{importante}
El vector de estado y los elementos orbitales son dos maneras equivalentes de describir la misma órbita kepleriana. Uno privilegia la formulación dinámica instantánea; el otro, la interpretación geométrica global.
\end{importante}

\section{Órbita osculatriz y elementos osculantes}
En un sistema real sometido a perturbaciones no keplerianas —por ejemplo, achatamiento terrestre, presión de radiación o terceros cuerpos— la trayectoria no es una cónica exacta. Sin embargo, en cualquier instante \(t_0\), puede construirse la \textbf{órbita osculatriz}: la cónica kepleriana que sería seguida por el cuerpo si, a partir de \(t_0\), desaparecieran todas las perturbaciones.

Sus propiedades principales son:
\begin{itemize}
    \item la órbita osculatriz es \textbf{tangente} a la trayectoria real en \(t_0\),
    \item comparte el mismo vector de estado \(\{\vec{r}(t_0),\vec{v}(t_0)\}\),
    \item sus elementos \((p(t), e(t), I(t), \Omega(t), \omega(t), f(t))\) dependen del tiempo y se denominan \textbf{elementos osculantes}.
\end{itemize}

La evolución temporal de estos elementos se describe mediante ecuaciones de variación de parámetros, como las ecuaciones de Lagrange o de Gauss, que relacionan las fuerzas perturbadoras con las derivadas temporales de los elementos orbitales.

\section{Síntesis final}
El desarrollo de esta clase cierra el puente entre la integración analítica del problema de dos cuerpos y su aplicación observacional:
\begin{enumerate}
    \item La geometría de la cónica se parametriza eficientemente mediante seis elementos orbitales.
    \item La orientación en el espacio se resuelve mediante tres rotaciones de Euler, cuya matriz conecta el sistema perifocal con el sistema del observador.
    \item La ambigüedad trigonométrica se resuelve mediante correcciones de cuadrante o mediante funciones trigonométricas con información de signo, como \(\operatorname{atan2}\).
    \item El vector de estado \(\{\vec{r},\vec{v}\}\) es suficiente y necesario para recuperar unívocamente los elementos orbitales.
    \item El concepto de órbita osculatriz permite extender el formalismo kepleriano a sistemas perturbados, sentando las bases para la teoría de perturbaciones y la propagación numérica de órbitas.
\end{enumerate}

\begin{importante}
Este marco constituye una de las columnas vertebrales de la astrodinámica moderna: conecta geometría, dinámica, observación y cómputo en un único lenguaje orbital coherente.
\end{importante}
\chapter{Clase 21: Determinación de órbita en el problema de dos cuerpos}

En el marco de la mecánica celeste newtoniana, el problema de dos cuerpos constituye la base analítica sobre la cual se construyen aproximaciones más complejas y se describe la dinámica de sistemas astronómicos reales. Hasta este punto del curso hemos trabajado principalmente en el problema directo: dada una configuración orbital completamente especificada, es posible predecir la posición y velocidad de los cuerpos en cualquier instante futuro. En esta clase invertimos el flujo de información: partiendo de un estado cinemático instantáneo conocido, aprenderemos a reconstruir la geometría completa de la trayectoria.

Este proceso se conoce como \textbf{determinación de órbita} y es fundamental tanto en astronomía observacional —donde se miden posiciones y velocidades de cuerpos recién descubiertos— como en ingeniería aeroespacial —donde se requiere navegación y control de satélites—.

El procedimiento descrito aquí se fundamenta en los principios de conservación del momento angular y de la energía, así como en la existencia del vector de excentricidad.

\section{Conjuntos de elementos orbitales}
La descripción geométrica de una cónica en el espacio tridimensional requiere exactamente seis parámetros independientes. Dependiendo del contexto físico y de las características dinámicas del cuerpo en estudio, se utilizan distintas parametrizaciones equivalentes.

\subsection{Elementos orbitales clásicos}
Se expresan como la tupla
\[
\vec{\Theta}_{\mathrm{cl}} = (p, e, I, \Omega, \omega, f).
\]
Sus significados son:
\begin{itemize}
    \item \(p\): semilatus rectum, que determina la escala de la cónica,
    \item \(e\): excentricidad, que fija la forma geométrica,
    \item \(I\): inclinación orbital, ángulo entre el plano orbital y el plano de referencia,
    \item \(\Omega\): longitud del nodo ascendente, orientación de la línea de nodos,
    \item \(\omega\): argumento del periapsis, orientación del eje mayor dentro del plano orbital,
    \item \(f\): anomalía verdadera, posición instantánea del cuerpo sobre la trayectoria.
\end{itemize}

\subsection{Elementos asteroidales}
Se usan convencionalmente para cuerpos con órbitas elípticas ligadas:
\[
\vec{\Theta}_{\mathrm{ast}} = (a, e, I, \Omega, \omega, f).
\]
Aquí el semilatus rectum \(p\) se reemplaza por el semieje mayor \(a\), relacionado mediante
\[
 p = a(1-e^2).
\]
Esta parametrización es especialmente útil porque \(a\) está ligado directamente a la energía específica relativa y al período orbital.

\subsection{Elementos cometarios}
Se emplean para cuerpos con trayectorias casi parabólicas o hiperbólicas:
\[
\vec{\Theta}_{\mathrm{com}} = (q, e, I, \Omega, \omega, f).
\]
En este caso se prefiere usar la distancia al periapsis \(q\), definida por
\[
 q = a(1-e) = \frac{p}{1+e}.
\]

La elección entre estos conjuntos es convencional y no altera la descripción física del sistema. Todas las representaciones son intercambiables mediante transformaciones algebraicas elementales.

\begin{importante}
Los elementos orbitales no son cantidades arbitrarias: constituyen una manera geométrica compacta de codificar exactamente la misma información que contiene el vector de estado \(\{\vec{r},\vec{v}\}\).
\end{importante}

\section{Relación entre el vector de estado y los elementos orbitales}
El estado completo de un sistema de dos cuerpos en un instante \(t\) se describe mediante el \textbf{vector de estado}:
\[
\vec{X}(t) =
\begin{pmatrix}
 x(t) \\
 y(t) \\
 z(t) \\
 v_x(t) \\
 v_y(t) \\
 v_z(t)
\end{pmatrix}
=
\begin{pmatrix}
 \vec{r}(t) \\
 \vec{v}(t)
\end{pmatrix}
\]
Este vector contiene seis componentes escalares que especifican posición y velocidad en un sistema de referencia inercial.

\subsection{¿Por qué seis elementos producen solo tres coordenadas espaciales?}
Si únicamente se dispone de la posición \(\vec{r} = (x,y,z)\), es imposible determinar la orientación espacial de la órbita ni su forma geométrica. La posición solo indica un punto por el que pasa la trayectoria, pero no revela hacia dónde se dirige el cuerpo, ni el plano en el cual se mueve, ni la curvatura local de la trayectoria.

Por esta razón, la proyección
\[
(p,e,I,\Omega,\omega,f) \longrightarrow (x,y,z)
\]
no es suficiente para reconstruir la dinámica completa.

\subsection{El difeomorfismo estado-elementos}
Cuando se incorpora también la velocidad \(\vec{v} = (v_x,v_y,v_z)\), el sistema de ecuaciones que relaciona el vector de estado con los elementos orbitales se vuelve completo y no degenerado. Matemáticamente, existe una correspondencia uno a uno entre el espacio de estados y el espacio de elementos orbitales, salvo en singularidades geométricas particulares, por ejemplo:
\begin{itemize}
    \item inclinación nula,
    \item excentricidad exactamente cero,
    \item casos límite parabólicos o rectilíneos.
\end{itemize}

Fuera de estas singularidades, la relación
\[
\vec{X} \longleftrightarrow (p,e,I,\Omega,\omega,f)
\]
es invertible y suave.

\section{El problema de determinación de órbita}
El núcleo de esta clase consiste en resolver analíticamente la transformación inversa: dado
\[
\vec{X} = (\vec{r},\vec{v})
\]
medido en un instante arbitrario, encontrar los seis elementos orbitales que describen la cónica kepleriana compatible con esas condiciones iniciales.

El procedimiento se basa en la construcción de tres vectores geométricos fundamentales derivados directamente del estado instantáneo y constantes en el problema no perturbado:
\begin{enumerate}
    \item \textbf{Momento angular específico relativo:}
\[
\vec{h} = \vec{r} \times \vec{v}
\]
    Este vector es perpendicular al plano orbital y su magnitud determina la velocidad areal.
    \item \textbf{Vector de excentricidad:}
\[
\vec{e} = \frac{\vec{v} \times \vec{h}}{\mu} - \frac{\vec{r}}{r}
\]
    Este vector yace en el plano orbital, apunta hacia el periapsis y cumple \(e = \|\vec{e}\|\).
    \item \textbf{Vector nodal:}
\[
\vec{n} = \hat{k} \times \vec{h}
\]
    donde \(\hat{k}\) es el vector unitario del eje \(z\). Este vector apunta a lo largo de la línea de nodos, específicamente hacia el nodo ascendente.
\end{enumerate}

Con estos tres vectores, la geometría completa de la órbita queda algebraicamente determinada.

\section{Construcción geométrica paso a paso}
A continuación se detalla el algoritmo analítico para extraer cada elemento orbital a partir de \(\vec{r}\) y \(\vec{v}\).

\subsection{Parámetros de tamaño y forma}
\[
 h = \|\vec{h}\|, \qquad p = \frac{h^2}{\mu}, \qquad e = \|\vec{e}\|
\]
Si se requiere el semieje mayor, por ejemplo en el caso elíptico o hiperbólico, entonces
\[
 a = \frac{p}{1-e^2}.
\]

\subsection{Inclinación orbital \texorpdfstring{$I$}{I}}
La inclinación es el ángulo entre \(\vec{h}\) y el eje \(z\) inercial:
\[
 \cos I = \frac{\vec{h} \cdot \hat{k}}{h} = \frac{h_z}{h}.
\]
Dado que \(I \in [0,\pi]\), la función \(\arccos\) produce aquí un valor único y correcto sin ambigüedad cuadrantal.

\subsection{Longitud del nodo ascendente \texorpdfstring{$\Omega$}{Omega}}
Se calcula a partir del vector nodal:
\[
 \cos\Omega = \frac{\vec{n} \cdot \hat{i}}{n} = \frac{n_x}{n}.
\]
Aquí aparece la primera ambigüedad de cuadrante, pues \(\arccos\) solo devuelve valores en \([0,\pi]\).

\subsection{Argumento del periapsis \texorpdfstring{$\omega$}{omega}}
Es el ángulo entre \(\vec{n}\) y \(\vec{e}\), medido dentro del plano orbital:
\[
 \cos\omega = \frac{\vec{n} \cdot \vec{e}}{ne}.
\]
También aquí el valor principal de \(\arccos\) no basta para decidir el cuadrante correcto.

\subsection{Anomalía verdadera \texorpdfstring{$f$}{f}}
Es el ángulo entre \(\vec{e}\) y \(\vec{r}\):
\[
 \cos f = \frac{\vec{e} \cdot \vec{r}}{er}.
\]
De nuevo, hace falta un criterio adicional para resolver la ambigüedad angular.

\section{La pesadilla de los cuadrantes}
Las funciones trigonométricas inversas estándar, en particular el arco coseno, son multivaluadas y sus implementaciones numéricas devuelven únicamente el valor principal en \([0,\pi]\). En mecánica celeste, los ángulos \(\Omega\), \(\omega\) y \(f\) deben definirse en el intervalo \([0,2\pi)\) para preservar la orientación correcta y el sentido del movimiento.

Resolver esta ambigüedad exige inspeccionar signos de proyecciones auxiliares.

\subsection{Corrección para \texorpdfstring{$\Omega$}{Omega}}
El nodo ascendente se identifica mediante el signo de la componente \(y\) del vector nodal:
\[
 \Omega =
\begin{cases}
\arccos\left(\frac{n_x}{n}\right), & \text{si } n_y \ge 0 \\
2\pi - \arccos\left(\frac{n_x}{n}\right), & \text{si } n_y < 0
\end{cases}
\]

\subsection{Corrección para \texorpdfstring{$\omega$}{omega}}
El argumento del periapsis requiere examinar la componente \(z\) del vector de excentricidad:
\[
 \omega =
\begin{cases}
\arccos\left(\frac{\vec{n} \cdot \vec{e}}{ne}\right), & \text{si } e_z \ge 0 \\
2\pi - \arccos\left(\frac{\vec{n} \cdot \vec{e}}{ne}\right), & \text{si } e_z < 0
\end{cases}
\]

\subsection{Corrección para \texorpdfstring{$f$}{f}}
La anomalía verdadera se corrige usando la velocidad radial
\[
 v_r = \frac{\vec{r} \cdot \vec{v}}{r}.
\]
Entonces,
\[
 f =
\begin{cases}
\arccos\left(\frac{\vec{e} \cdot \vec{r}}{er}\right), & \text{si } v_r \ge 0 \\
2\pi - \arccos\left(\frac{\vec{e} \cdot \vec{r}}{er}\right), & \text{si } v_r < 0
\end{cases}
\]

Estas correcciones garantizan que los elementos orbitales recuperados sean geométricamente consistentes con el vector de estado original, evitando reflexiones espurias introducidas por el uso ingenuo de las funciones trigonométricas inversas.

\begin{importante}
Toda implementación seria de determinación de órbita debe resolver explícitamente la ambigüedad cuadrantal. De lo contrario, puede obtener una órbita geométricamente correcta en forma pero incorrecta en orientación o en sentido de movimiento.
\end{importante}

\section{Órbita osculatriz y elementos osculatrices}
La determinación de órbita descrita hasta aquí no requiere que la trayectoria real sea una cónica exacta. En presencia de perturbaciones —interacciones con terceros cuerpos, arrastre atmosférico, presión de radiación, achatamiento terrestre, entre otras— la dinámica se desvía del modelo kepleriano ideal. Sin embargo, en cualquier instante \(t\), el vector de estado \(\vec{X}(t)\) sigue estando bien definido.

Al aplicar el algoritmo de determinación de órbita a \(\vec{X}(t)\) en un instante arbitrario, obtenemos un conjunto de elementos \(\vec{\Theta}(t)\) que describen una cónica hipotética que:
\begin{enumerate}
    \item tiene como foco el origen de coordenadas,
    \item pasa exactamente por \(\vec{r}(t)\),
    \item es tangente a la trayectoria real en ese punto porque comparte el mismo vector velocidad.
\end{enumerate}

A esta cónica se le denomina \textbf{órbita osculatriz}. Sus parámetros se conocen como \textbf{elementos osculatrices}. En el problema puro de dos cuerpos, estos elementos son constantes en el tiempo. En sistemas perturbados, se convierten en funciones del tiempo
\[
\vec{\Theta}(t),
\]
cuya evolución está gobernada por ecuaciones de variación de parámetros, como las ecuaciones de Lagrange o de Gauss.

\section{Síntesis analítica y comentarios finales}
El proceso de determinación de órbita resume elegantemente la dualidad geométrica y dinámica de la mecánica celeste:
\begin{itemize}
    \item la \textbf{cinemática}, a través de \((\vec{r},\vec{v})\), proporciona el estado instantáneo,
    \item la \textbf{dinámica}, a través de la conservación de \(\vec{h}\) y \(\vec{e}\), revela la estructura invariante del movimiento,
    \item la \textbf{geometría}, a través de \((I,\Omega,\omega,f)\), traduce estas invariantes en parámetros observables y predictivos.
\end{itemize}

Es importante recordar que este formalismo supone un sistema de referencia inercial bien definido. Las transformaciones entre distintos marcos requieren matrices de rotación basadas en ángulos de Euler, tal como se discutió en la clase anterior.

\begin{importante}
La determinación de órbita constituye uno de los puentes más poderosos entre teoría y observación: transforma un estado instantáneo medido en una trayectoria completa, interpretable y predecible dentro del formalismo kepleriano.
\end{importante}
\chapter{Clase 23: La velocidad relativa y el hodógrafo en el problema de dos cuerpos}

En el desarrollo de la mecánica celeste newtoniana, la descripción completa del movimiento relativo de un sistema de dos cuerpos requiere conocer no solo la geometría de la trayectoria, sino también la cinemática instantánea del cuerpo en cada punto de la misma. En clases anteriores se estableció cómo, dados los elementos orbitales clásicos \((p,e,I,\Omega,\omega,f)\), es posible determinar el vector de posición \(\vec{r}\) en cualquier instante. En esta clase abordaremos el problema complementario: cómo calcular el vector velocidad \(\vec{v}\), qué propiedades geométricas posee y por qué su comportamiento revela una de las estructuras más elegantes de toda la mecánica celeste: el \textbf{hodógrafo}.

El desarrollo se apoya en la relación entre el vector de estado y los elementos orbitales, en la conservación de la energía específica y del momento angular específico, y en la estructura del vector de excentricidad.

\section{Del vector de estado a los elementos orbitales}
El estado dinámico completo de un sistema de dos cuerpos en un instante arbitrario \(t\) se describe mediante el \textbf{vector de estado}:
\[
\vec{X}(t) =
\begin{pmatrix}
 x(t) \\
 y(t) \\
 z(t) \\
 v_x(t) \\
 v_y(t) \\
 v_z(t)
\end{pmatrix}
=
\begin{pmatrix}
 \vec{r}(t) \\
 \vec{v}(t)
\end{pmatrix}
\]
Este vector contiene seis componentes escalares que especifican posición y velocidad en un sistema de referencia inercial.

\subsection{¿Por qué seis elementos orbitales producen solo tres coordenadas espaciales?}
Si únicamente se dispone del vector posición \(\vec{r}=(x,y,z)\), es imposible determinar la orientación espacial de la órbita ni su curvatura. La posición solo indica un punto geométrico por el que pasa la trayectoria, pero no revela la dirección del movimiento, el plano orbital ni si la curva es cerrada o abierta. Matemáticamente, la aplicación
\[
(p,e,I,\Omega,\omega,f) \longrightarrow (x,y,z)
\]
no es invertible, pues proyecta seis parámetros sobre solo tres coordenadas y destruye información esencial sobre la tangente a la curva.

\subsection{El difeomorfismo estado-elementos}
Cuando se incorpora también la velocidad \(\vec{v}=(v_x,v_y,v_z)\), el sistema de ecuaciones que relaciona el estado cinemático con los elementos orbitales se vuelve completo y no degenerado. Existe entonces una correspondencia biyectiva, suave y con jacobiano no nulo entre el espacio de estados y el espacio de elementos orbitales, salvo en singularidades geométricas especiales. En forma simbólica:
\[
\vec{X} \longleftrightarrow (p,e,I,\Omega,\omega,f)
\]
La velocidad completa la información geométrica porque fija el plano orbital mediante \(\vec{h}=\vec{r}\times\vec{v}\), determina la excentricidad y el tamaño orbital a través de la energía específica, y localiza el periapsis mediante el vector de excentricidad.

\begin{importante}
Conocer \((\vec{r},\vec{v})\) es equivalente a conocer la cónica completa y la posición instantánea del cuerpo sobre ella. Esa equivalencia es la base analítica de toda determinación de órbita.
\end{importante}

\section{El problema de la velocidad relativa}
Una vez determinada la geometría orbital, surge una pregunta natural: ¿cuál es la magnitud y dirección del vector velocidad relativa en un punto arbitrario de la trayectoria? Resolver este problema no requiere integrar nuevamente las ecuaciones de movimiento; basta con explotar las constantes de movimiento ya conocidas: la energía específica relativa \(\epsilon\) y el momento angular específico relativo \(h\).

\subsection{Energía específica relativa en el periapsis}
En el periapsis, el vector posición y el vector velocidad son perpendiculares. Por tanto, la magnitud del momento angular específico se simplifica a
\[
 h = rv.
\]
Además, la distancia al periapsis es
\[
 q = a(1-e).
\]
Sustituyendo en la definición de energía específica:
\[
\epsilon = \frac{v^2}{2} - \frac{\mu}{r} = \frac{h^2}{2q^2} - \frac{\mu}{q}.
\]
Usando la relación geométrica
\[
 h^2 = \mu p = \mu a(1-e^2)
\]
y que \(q=a(1-e)\), se obtiene, tras una simplificación algebraica directa,
\[
\epsilon = -\frac{\mu}{2a}.
\]
Este resultado es fundamental: la energía específica relativa de una órbita kepleriana depende exclusivamente del semieje mayor \(a\), independientemente de la excentricidad.

\subsection{Energía específica en un punto arbitrario}
Dado que \(\epsilon\) es una constante de movimiento, su valor en cualquier punto de la trayectoria coincide con el obtenido en el periapsis. Igualando la expresión general con el valor anterior:
\[
\frac{v^2}{2} - \frac{\mu}{r} = -\frac{\mu}{2a}
\]
Despejando \(v^2\), obtenemos inmediatamente la ecuación fundamental de la rapidez orbital.

\section{Magnitud de la velocidad: la ecuación vis-viva}
La manipulación algebraica de la expresión anterior conduce directamente a
\[
 v^2 = \mu\left(\frac{2}{r} - \frac{1}{a}\right).
\]
A esta relación se le conoce como la \textbf{ecuación vis-viva}.

\subsection{Propiedades analíticas de la vis-viva}
\begin{enumerate}
    \item \textbf{Independencia direccional:} dada una distancia \(r\) y una rapidez \(v\), el semieje mayor queda determinado unívocamente por
\[
 a = \frac{\mu}{\frac{2\mu}{r} - v^2}.
\]
    Nótese que \(a\) no depende de la dirección del vector velocidad. Dos cuerpos lanzados desde la misma posición con la misma rapidez pero direcciones distintas tienen el mismo tamaño orbital, aunque diferente orientación y excentricidad.
    \item \textbf{Límites físicos:} para órbitas elípticas \((a>0)\), se requiere \(v < \sqrt{2\mu/r}\), que es la velocidad de escape local. Para trayectorias parabólicas, \(a \to \infty\) y \(v=\sqrt{2\mu/r}\). Para órbitas hiperbólicas \((a<0)\), se cumple \(v > \sqrt{2\mu/r}\).
    \item \textbf{Velocidades características:} en una elipse, la velocidad es máxima en el periapsis y mínima en el apoapsis:
\[
 v_p = \sqrt{\frac{\mu}{a}\frac{1+e}{1-e}}, \qquad
 v_a = \sqrt{\frac{\mu}{a}\frac{1-e}{1+e}}.
\]
\end{enumerate}

\begin{importante}
La ecuación vis-viva conecta de manera directa la cinemática local con la energía global de la órbita. Es una de las fórmulas más útiles de toda la astrodinámica.
\end{importante}

\section{Dirección de la velocidad y ángulo de trayectoria}
Conocer la magnitud \(v=\|\vec{v}\|\) no basta para determinar el vector completo. La dirección se obtiene aprovechando la definición del momento angular específico:
\[
\vec{h} = \vec{r} \times \vec{v}.
\]
Tomando magnitudes:
\[
 h = rv\sin\phi,
\]
donde \(\phi\) es el ángulo entre \(\vec{r}\) y \(\vec{v}\), conocido como \textbf{ángulo de trayectoria}.

Despejando:
\[
 \sin\phi = \frac{h}{rv}.
\]
Esta relación permite determinar la inclinación del vector velocidad respecto a la dirección local horizontal.

\subsection{Ambigüedad cuadrantal y velocidad radial}
Como \(h\), \(r\) y \(v\) son positivos, el valor de \(\sin\phi\) restringe \(\phi\) al primer o segundo cuadrante. Para resolver esta ambigüedad se proyecta la velocidad sobre la dirección radial:
\[
 v_r = \vec{v} \cdot \hat{r} = v\cos\phi.
\]
Alternativamente, derivando la ecuación de la cónica, se obtiene
\[
 v_r = \frac{\mu}{h}e\sin f.
\]
El signo de \(v_r\), o equivalentemente de \(\sin f\), determina si el cuerpo se aleja del periapsis \((0<f<\pi)\) o se acerca a él \((\pi<f<2\pi)\), fijando así el cuadrante correcto de \(\phi\).

\section{El hodógrafo: geometría en el espacio de velocidades}
Uno de los resultados más sorprendentes del problema de dos cuerpos es que el vector velocidad describe una circunferencia en el espacio de velocidades. A este lugar geométrico se le denomina \textbf{hodógrafo}.

\subsection{Derivación analítica}
Partimos de la definición del vector de excentricidad:
\[
\vec{e} = \frac{\vec{v} \times \vec{h}}{\mu} - \hat{r}.
\]
Evaluando las componentes de \(\vec{v}\) en el sistema perifocal, donde el eje \(x\) apunta al periapsis y el eje \(y\) es transversal, se obtiene
\[
 v_x = -\frac{\mu}{h}\sin f, \qquad
 v_y = \frac{\mu}{h}(e+\cos f), \qquad
 v_z = 0.
\]
Eliminando la anomalía verdadera \(f\) mediante \(\sin^2 f + \cos^2 f = 1\), resulta
\[
 v_x^2 + \left(v_y - \frac{\mu e}{h}\right)^2 = \left(\frac{\mu}{h}\right)^2.
\]
Esta es la ecuación canónica de una circunferencia en el plano \(v_xv_y\).

\subsection{Interpretación geométrica del hodógrafo}
La circunferencia del hodógrafo tiene:
\begin{itemize}
    \item \textbf{Centro:} \(\left(0,\frac{\mu e}{h}\right)\),
    \item \textbf{Radio:} \(R_v = \frac{\mu}{h}\).
\end{itemize}

Su comportamiento en casos particulares es muy instructivo:
\begin{itemize}
    \item si \(e=0\), el centro coincide con el origen y la velocidad es constante en magnitud,
    \item si \(e=1\), la circunferencia pasa por el origen, reflejando que la velocidad tiende a cero en el infinito parabólico,
    \item si \(e>1\), el origen queda fuera de la circunferencia, indicando que la velocidad nunca se anula.
\end{itemize}

El hodógrafo transforma así un problema de geometría cónica en el espacio de posiciones en un problema circular en el espacio de velocidades.

\begin{importante}
El hodógrafo es una de las manifestaciones más elegantes de la simetría oculta del problema kepleriano: la complejidad geométrica de las cónicas se reduce a una circunferencia en el espacio de velocidades.
\end{importante}

\section{Componentes de la velocidad en el sistema inercial}
Hasta ahora hemos trabajado en el sistema perifocal, natural de la cónica. Para predecir el vector de estado en un marco inercial arbitrario, es necesario rotar las componentes de la velocidad mediante los ángulos de Euler \((\Omega,I,\omega)\). La matriz
\[
\mathbf{M}(\Omega,I,\omega)=\mathbf{R}_z(\Omega)\mathbf{R}_x(I)\mathbf{R}_z(\omega)
\]
transforma vectores desde el sistema perifocal al sistema inercial.

Aplicando esta transformación a las componentes perifocales \((v_x,v_y,0)\), se obtienen las componentes inerciales explícitas:
\[
\dot{x} = \frac{\mu}{h}\left[-\cos\Omega\sin(\omega+f)-\cos I\sin\Omega\cos(\omega+f)\right]-\frac{\mu e}{h}\left(\cos\Omega\sin\omega+\cos\omega\cos I\sin\Omega\right)
\]
\[
\dot{y} = \frac{\mu}{h}\left[-\sin\Omega\sin(\omega+f)+\cos I\cos\Omega\cos(\omega+f)\right]+\frac{\mu e}{h}\left(-\sin\Omega\sin\omega+\cos\omega\cos I\cos\Omega\right)
\]
\[
\dot{z} = \frac{\mu}{h}\left[\sin I\cos(\omega+f)+e\cos\omega\sin I\right].
\]

Estas expresiones cierran el ciclo de predicción: dados los elementos orbitales y un valor de \(f\), es posible reconstruir completamente el vector de estado \(\vec{X}=(\vec{r},\vec{v})\) en cualquier marco de referencia inercial.

\section{Síntesis analítica y comentarios finales}
El estudio de la velocidad en el problema de dos cuerpos revela una conexión profunda entre energía, geometría y álgebra vectorial:
\begin{enumerate}
    \item La ecuación vis-viva vincula directamente la cinemática local con el tamaño global de la órbita, mostrando que la energía específica dicta el destino energético del sistema.
    \item El hodógrafo muestra que la dinámica kepleriana posee una simetría oculta en el espacio de velocidades, donde la complejidad de las cónicas se reduce a movimientos circulares simples.
    \item La transformación inercial de la velocidad confirma que, aunque la magnitud y dirección relativa dependen únicamente de \(r\) y \(f\), su proyección espacial exige la rotación completa mediante los ángulos orbitales.
\end{enumerate}

\begin{importante}
La velocidad relativa no es solo una derivada temporal de la posición: es un objeto geométrico con estructura propia, capaz de revelar invariantes profundas del problema kepleriano y de conectar la teoría analítica con las aplicaciones prácticas de la astrodinámica.
\end{importante}
\chapter{Clase 24: La órbita en el tiempo (parte 1)}

En las clases precedentes se estableció un puente bidireccional entre el vector de estado instantáneo
\[
\vec{X} = (\vec{r},\vec{v})
\]
y los elementos orbitales clásicos. Aprendimos que, conocidos los seis elementos orbitales clásicos —\(a\), \(e\), \(I\), \(\Omega\), \(\omega\) y \(f\)—, es posible reconstruir analíticamente la posición espacial \(\vec{r}\) y el vector velocidad \(\vec{v}\) del sistema relativo mediante rotaciones de Euler y ecuaciones paramétricas planas. Sin embargo, esta descripción geométrica tiene una limitación fundamental: no especifica \textbf{cómo} evoluciona el cuerpo a lo largo de su trayectoria con el paso del tiempo.

De los seis elementos orbitales, cinco —\(a\), \(e\), \(I\), \(\Omega\) y \(\omega\)— son constantes de movimiento en el problema no perturbado de dos cuerpos. El único que varía con el tiempo es la anomalía verdadera \(f\). La pregunta central de esta clase es, por tanto: ¿cómo calcular la anomalía verdadera como función del tiempo, \(f(t)\)? Este desafío constituye el núcleo del problema temporal de Kepler.

La mecánica newtoniana indica que, entre todas las cantidades cinemáticas y geométricas asociadas a la cónica, el área barrida por el vector relativo es la única que posee una dependencia temporal explícita, lineal y completamente predecible.

\section{La ley de las áreas y la predicción temporal}
La constancia del momento angular específico relativo
\[
\vec{h} = \vec{r} \times \vec{v}
\]
implica directamente que la tasa de barrido del área por el vector posición es invariable. Matemáticamente, el área diferencial \(dA\) barrida en un intervalo \(dt\) corresponde a la mitad del área del paralelogramo formado por \(\vec{r}\) y \(d\vec{r}\):
\[
 dA = \frac{1}{2}\,\|\vec{r} \times d\vec{r}\| = \frac{1}{2}\,\|\vec{r} \times \vec{v}\,dt\| = \frac{1}{2}h\,dt
\]

Integrando esta relación desde el instante de paso por el periapsis \(t_p\), donde por definición el área barrida es nula, hasta un tiempo arbitrario \(t\), obtenemos:
\[
 A(t) = \frac{1}{2}h(t-t_p).
\]

Esta expresión revela que el área crece linealmente con el tiempo. El problema analítico consiste ahora en invertir esta relación: dado un tiempo \(t\), conocemos \(A(t)\), pero necesitamos traducir ese valor escalar en una posición angular sobre la cónica.

La anomalía verdadera \(f\) no es adecuada para esta integración directa debido a la complejidad de la ecuación polar
\[
 r = \frac{p}{1+e\cos f}.
\]
Por esta razón, se introduce un parámetro angular auxiliar que simplifica notablemente la geometría del problema: la \textbf{anomalía excéntrica} \(E\).

\begin{importante}
El problema del tiempo en órbitas keplerianas no se resuelve directamente con la anomalía verdadera. La clave es encontrar una variable intermedia cuya relación con el área sea algebraicamente simple.
\end{importante}

\section{La anomalía excéntrica y su construcción geométrica}
Para definir \(E\), consideremos una elipse con semieje mayor \(a\) y semieje menor \(b\), centrada en el origen de un sistema cartesiano auxiliar \((x_c,y_c)\). Si circunscribimos una circunferencia de radio \(a\) alrededor de la elipse, cualquier punto \(P\) sobre la curva elíptica puede proyectarse paralelamente al eje menor hasta intersectar la circunferencia en un punto \(Q\). La anomalía excéntrica \(E\) se define como el ángulo polar de \(Q\), medido desde el semieje mayor positivo y con centro en el centro geométrico de la elipse.

Esta construcción geométrica permite escribir las coordenadas del punto sobre la elipse en forma paramétrica:
\[
 x_c = a\cos E,
\]
\[
 y_c = b\sin E.
\]

A diferencia de la anomalía verdadera, que se mide desde el foco y cuya velocidad angular varía según la segunda ley de Kepler, la anomalía excéntrica no tiene una interpretación cinemática directa. Su utilidad reside en sus propiedades trigonométricas y algebraicas, que la hacen ideal para la integración temporal.

\section{Relación entre la anomalía verdadera y la anomalía excéntrica}
Es posible expresar la distancia radial \(r\) directamente en función de \(E\) sin pasar primero por la anomalía verdadera. Trasladando el origen desde el centro geométrico hasta el foco, es decir, desplazando una distancia \(ae\) en la dirección del periapsis, y aplicando el teorema de Pitágoras, se obtiene:
\[
 r^2 = (a\cos E - ae)^2 + (b\sin E)^2.
\]

Usando la relación
\[
 b^2 = a^2(1-e^2)
\]
y desarrollando algebraicamente, los términos se simplifican hasta dar la expresión compacta
\[
 r = a(1-e\cos E).
\]

Esta fórmula es mucho más manejable analíticamente que la ecuación polar clásica y evita la singularidad trigonométrica del denominador.

Para conectar \(E\) con la anomalía verdadera \(f\), aplicamos identidades de ángulo medio a la geometría focal. El resultado exacto es
\[
 \tan\left(\frac{f}{2}\right) = \sqrt{\frac{1+e}{1-e}}\,\tan\left(\frac{E}{2}\right).
\]

De aquí se deduce que, para órbitas elípticas \((0<e<1)\), el valor absoluto de la anomalía verdadera suele ser mayor que el de la anomalía excéntrica, salvo en los puntos singulares del periapsis y del apoapsis, donde ambas coinciden. Esto refleja físicamente que el cuerpo barre más ángulo verdadero cerca del periapsis, donde su velocidad orbital es mayor.

\section{Cálculo del área del sector elíptico}
El paso crucial para vincular la geometría con el tiempo consiste en calcular analíticamente el área \(\Delta A\) del sector elíptico barrido desde el periapsis hasta la posición angular correspondiente a \(E\).

Utilizando la propiedad de afinidad entre la elipse y su circunferencia circunscrita, el área del sector elíptico se obtiene escalando el área del sector circular correspondiente y restando el área del triángulo formado por el foco, el centro y el punto proyectado. El resultado final puede escribirse como
\[
 \Delta A = \frac{1}{2}ab(E-e\sin E).
\]

Cuando \(E = 2\pi\), recuperamos el área total de la elipse:
\[
 A_{\text{elipse}} = \pi ab,
\]
lo que confirma la consistencia geométrica de la expresión anterior.

\begin{importante}
La fórmula \(\Delta A = \frac{1}{2}ab(E-e\sin E)\) es el verdadero punto de contacto entre la geometría elíptica y la dinámica temporal del problema de Kepler.
\end{importante}

\section{Deducción de la ecuación de Kepler}
Finalmente, igualamos la expresión dinámica del área barrida en función del tiempo con la expresión geométrica en función de la anomalía excéntrica:
\[
 \frac{1}{2}h(t-t_p) = \frac{1}{2}ab(E-e\sin E).
\]

Sustituimos las relaciones geométricas conocidas del problema de dos cuerpos:
\[
 b = a\sqrt{1-e^2},
\qquad
 h = \sqrt{\mu a(1-e^2)}.
\]

Al simplificar, el factor \(\sqrt{1-e^2}\) se cancela y la ecuación se reorganiza como
\[
 \sqrt{\frac{\mu}{a^3}}(t-t_p) = E - e\sin E.
\]

Definimos ahora dos cantidades fundamentales:
\begin{itemize}
    \item el \textbf{movimiento medio}
\[
 n \equiv \sqrt{\frac{\mu}{a^3}},
\]
    que representa la velocidad angular promedio que tendría el cuerpo si se moviera en una órbita circular del mismo período;
    \item la \textbf{anomalía media}
\[
 M \equiv n(t-t_p),
\]
    una variable angular auxiliar que crece linealmente con el tiempo.
\end{itemize}

Con estas definiciones obtenemos la ecuación fundamental de la propagación temporal orbital:
\[
 M = E - e\sin E.
\]

Esta es la \textbf{ecuación de Kepler}. Se trata de una ecuación trascendental que relaciona linealmente el tiempo, a través de \(M\), con la posición geométrica, a través de \(E\). Dado un tiempo \(t\), calcular \(M\) es inmediato; sin embargo, despejar \(E\) requiere métodos numéricos o aproximaciones analíticas, ya que \(E\) aparece tanto linealmente como dentro de una función trigonométrica.

\section{Síntesis analítica y comentarios finales}
En esta clase se estableció la estructura matemática necesaria para resolver el problema de la evolución temporal en el sistema de dos cuerpos. Los resultados clave son:
\begin{enumerate}
    \item El área barrida es la única cantidad que evoluciona linealmente con el tiempo debido a la conservación de \(\vec{h}\).
    \item La anomalía excéntrica \(E\) actúa como un parámetro geométrico intermedio que linealiza la dependencia de la distancia radial
\[
 r = a(1-e\cos E)
\]
y simplifica el cálculo del área sectorial.
    \item La ecuación de Kepler
\[
 M = E - e\sin E
\]
cierra el ciclo entre tiempo y geometría, permitiendo predecir la posición orbital para cualquier \(t>t_p\).
\end{enumerate}

La resolución práctica de esta ecuación para \(E\), seguida de la transformación a \(f\) y luego a \(\vec{r}(t)\), constituye la base de todos los algoritmos de propagación orbital en astrodinámica y mecánica celeste.

\begin{importante}
La dificultad del problema temporal no radica en la física, sino en la inversión analítica de la ecuación de Kepler. Todo el resto del formalismo orbital descansa sobre esa inversión.
\end{importante}
\chapter{Clase 26: La órbita en el tiempo (parte 2)}

En la clase anterior se establecieron las bases geométricas para relacionar el estado cinemático instantáneo con la forma de la trayectoria mediante los elementos orbitales. Sin embargo, una descripción puramente geométrica no especifica de manera explícita cómo evoluciona la posición del cuerpo a lo largo de su órbita en función del tiempo físico. En esta clase abordamos el núcleo del \textbf{problema de Kepler}: determinar la posición y la velocidad de un sistema de dos cuerpos en un instante arbitrario \(t\), partiendo de condiciones iniciales conocidas y de las constantes de movimiento.

El desarrollo se fundamenta en la integración del teorema de las áreas, la introducción de la anomalía media como parámetro temporal lineal, la deducción y resolución de la ecuación de Kepler y, finalmente, la síntesis de un algoritmo completo de propagación orbital. Además, se introducen las funciones de Lagrange \(f\) y \(g\), una formulación alternativa que permite propagar el estado sin recurrir explícitamente a los elementos orbitales.

\section{El teorema de áreas y la tercera ley planetaria}
\subsection{Integración del barrido areal}
Como se demostró previamente, la constancia del momento angular específico relativo
\[
\vec{h} = \vec{r} \times \vec{v}
\]
implica que la tasa de barrido del área es constante:
\[
\frac{dA}{dt} = \frac{1}{2}h.
\]

Integrando desde el instante de paso por el periapsis \(t_p\) hasta un tiempo arbitrario \(t\), el área acumulada es
\[
A(t) = \frac{1}{2}h(t-t_p).
\]

Esta relación es la única cantidad geométrica del problema que posee una dependencia temporal explícita, lineal y predecible.

\subsection{El período orbital y la tercera ley de Kepler}
Para una órbita cerrada, el cuerpo completa una revolución cuando el área barrida coincide con el área total de la elipse,
\[
A_{\text{elipse}} = \pi ab,
\]
donde \(a\) es el semieje mayor y \(b\) el semieje menor. Igualando con la expresión dinámica para un período completo \(T\):
\[
\pi ab = \frac{1}{2}hT.
\]

Sustituyendo las relaciones
\[
 b = a\sqrt{1-e^2},
 \qquad
 h = \sqrt{\mu a(1-e^2)},
\]
y simplificando, obtenemos
\[
 T = \frac{2\pi a^{3/2}}{\sqrt{\mu}}.
\]
Elevando al cuadrado:
\[
 T^2 = \frac{4\pi^2}{\mu}a^3.
\]

Esta es la formulación clásica de la \textbf{tercera ley de Kepler}: el cuadrado del período orbital es proporcional al cubo del semieje mayor, con una constante que depende solo del parámetro gravitacional del sistema.

\subsection{El movimiento medio y el teorema armónico}
Definimos la \textbf{velocidad angular media} o \textbf{movimiento medio} \(n\) como
\[
 n \equiv \frac{2\pi}{T} = \sqrt{\frac{\mu}{a^3}}.
\]
Esto conduce al \textbf{teorema armónico}:
\[
 n^2 a^3 = \mu.
\]

El movimiento medio es constante para una órbita kepleriana no perturbada y será la herramienta temporal básica para propagar la órbita.

\begin{importante}
Mientras la anomalía verdadera cambia de manera no uniforme, la anomalía media crece linealmente con el tiempo. Ese contraste es la clave para resolver el problema temporal kepleriano.
\end{importante}

\section{La anomalía media y la ecuación de Kepler}
\subsection{Definición de la anomalía media}
La anomalía media \(M\) es un parámetro angular auxiliar que crece linealmente con el tiempo a partir del paso por el periapsis:
\[
 M(t) = n(t-t_p).
\]

Geométricamente, \(M\) representa el ángulo que habría barrido un cuerpo ficticio si se moviera sobre la circunferencia auxiliar con velocidad angular constante \(n\). No corresponde a un ángulo físico real en el espacio, pero su linealidad temporal la convierte en la variable ideal para describir la evolución temporal de la órbita.

\subsection{Deducción de la ecuación de Kepler}
El vínculo entre el tiempo, a través de \(M\), y la geometría orbital, a través de \(E\), se obtiene igualando la expresión geométrica del área del sector elíptico con la expresión dinámica del área barrida. Del análisis geométrico se tiene:
\[
 \Delta A = \frac{1}{2}ab(E-e\sin E).
\]

Igualando con
\[
\Delta A = \frac{1}{2}h(t-t_p)
\]
y usando que
\[
 h = nab,
\]
se obtiene la \textbf{ecuación de Kepler}:
\[
 M = E - e\sin E.
\]

Esta es una ecuación trascendental que relaciona la anomalía media \(M\), conocida a partir del tiempo transcurrido, con la anomalía excéntrica \(E\), que determina la posición geométrica. Como \(E\) aparece tanto linealmente como dentro de una función trigonométrica, no puede despejarse mediante operaciones algebraicas elementales.

\subsection{Interpretación geométrica original}
Kepler interpretó esta ecuación comparando áreas en la circunferencia circunscrita a la elipse. La anomalía media es proporcional al área del sector circular menos el área de un triángulo correctivo. Esta construcción muestra por qué \(M\) y \(E\) difieren: la excentricidad \(e\) introduce una corrección geométrica que rompe la proporcionalidad lineal entre tiempo y ángulo real.

\section{Métodos de resolución de la ecuación de Kepler}
Dada la naturaleza trascendental de
\[
 M = E - e\sin E,
\]
su resolución requiere métodos numéricos o aproximaciones analíticas. Estos métodos son esenciales en astrodinámica moderna.

\subsection{Método de Newton--Raphson}
El método más eficiente y ampliamente usado consiste en plantear la búsqueda de raíces de la función
\[
 F(E) = E - e\sin E - M.
\]

Aplicando Newton--Raphson:
\[
 E_{k+1} = E_k - \frac{F(E_k)}{F'(E_k)},
\]
y dado que
\[
 F'(E) = 1 - e\cos E,
\]
obtenemos la recurrencia
\[
 E_{k+1} = E_k - \frac{E_k - e\sin E_k - M}{1 - e\cos E_k}.
\]

Este método converge cuadráticamente y suele requerir solo unas pocas iteraciones para alcanzar alta precisión.

\subsection{Método del punto fijo}
Kepler propuso originalmente una iteración más simple, reescribiendo la ecuación como
\[
 E = M + e\sin E.
\]
La recurrencia asociada es
\[
 E_{k+1} = M + e\sin E_k.
\]

Este método converge linealmente y es más lento que Newton--Raphson, sobre todo cuando \(e\) se aproxima a 1, pero tiene la ventaja de no involucrar derivadas ni denominadores sensibles.

\subsection{Solución analítica por series de Fourier}
Para aplicaciones teóricas y perturbativas, es útil expresar \(E\) como una serie de Fourier en términos de \(M\). Usando coeficientes de Bessel \(J_n\), puede escribirse:
\[
 E = M + \sum_{n=1}^{\infty} \frac{2}{n}J_n(ne)\sin(nM).
\]

Esta serie converge para \(e<1\). Para excentricidades pequeñas \((e \ll 1)\), truncando en orden \(e^2\), se obtiene la aproximación clásica
\[
 E \approx M + e\sin M + \frac{e^2}{2}\sin 2M.
\]

\subsection{Ejemplo práctico: posición de la Tierra}
Supongamos la posición de la Tierra 60 días después del paso por el periapsis. Si
\[
 n \approx 0.9856^\circ/\text{d},
\]
entonces la anomalía media es aproximadamente
\[
 M = 59.14^\circ.
\]
Para obtener la posición real se resuelve
\[
 59.14^\circ = E - 0.0167\sin E.
\]
Usando Newton--Raphson con \(E_0=M\), la convergencia es rápida y produce un valor muy próximo a \(E \approx 59.38^\circ\). Luego, usando la relación entre \(E\) y \(f\), se obtiene la anomalía verdadera y, finalmente, la posición física del planeta en su órbita.

\begin{importante}
El problema de Kepler no consiste en integrar nuevamente la ecuación de movimiento. Consiste en invertir de manera eficiente y precisa una ecuación trascendental para recuperar la geometría orbital a partir del tiempo.
\end{importante}

\section{Síntesis del algoritmo de predicción orbital}
El problema de los dos cuerpos en el tiempo se resuelve completamente mediante una secuencia algorítmica que transforma un estado inicial
\[
\vec{X}(t_0) = (\vec{r}_0,\vec{v}_0)
\]
en un estado futuro \(\vec{X}(t)\). Los pasos fundamentales son:
\begin{enumerate}
    \item Calcular los vectores invariantes:
\[
\vec{h} = \vec{r}_0 \times \vec{v}_0,
\qquad
\vec{e} = \frac{\vec{v}_0 \times \vec{h}}{\mu} - \frac{\vec{r}_0}{r_0}.
\]
    \item Determinar los elementos escalares:
\[
 p = \frac{h^2}{\mu}, \qquad e = \|\vec{e}\|, \qquad a = \frac{p}{1-e^2}, \qquad n = \sqrt{\frac{\mu}{|a|^3}}.
\]
    \item Calcular la anomalía inicial \(f_0\), luego \(E_0\), y después
\[
 M_0 = E_0 - e\sin E_0.
\]
    \item Propagar temporalmente:
\[
 M(t) = M_0 + n(t-t_0).
\]
    \item Resolver la ecuación de Kepler para obtener \(E(t)\).
    \item Convertir a anomalía verdadera:
\[
 f(t) = 2\arctan\left(\sqrt{\frac{1+e}{1-e}}\tan\frac{E}{2}\right).
\]
    \item Reconstruir \(\vec{r}(t)\) y \(\vec{v}(t)\) usando las ecuaciones paramétricas de la cónica y las rotaciones de Euler \((\Omega,i,\omega)\).
\end{enumerate}

Esta secuencia constituye la base de todos los algoritmos clásicos de propagación orbital.

\section{Las funciones de Lagrange \texorpdfstring{$f$}{f} y \texorpdfstring{$g$}{g}}
Existe una formulación alternativa para propagar el estado orbital que evita la conversión explícita a elementos orbitales y, por tanto, evita singularidades numéricas en casos de excentricidad o inclinación nula. Como el movimiento kepleriano es plano, el vector posición en cualquier instante puede expresarse como una combinación lineal del vector posición y velocidad iniciales:
\[
 \vec{r}(t) = f(t,t_0)\,\vec{r}_0 + g(t,t_0)\,\vec{v}_0.
\]

Derivando respecto al tiempo, la velocidad resulta
\[
 \vec{v}(t) = \dot{f}(t,t_0)\,\vec{r}_0 + \dot{g}(t,t_0)\,\vec{v}_0.
\]

Las funciones escalares \(f\) y \(g\), junto con sus derivadas, dependen únicamente de la geometría de la trayectoria y del tiempo transcurrido. En el marco de las \textbf{variables universales}, estas funciones se escriben de manera unificada para elipses, parábolas e hipérbolas usando las funciones de Stumpff:
\[
 f = 1 - \frac{\mu}{r_0}x^2 c_2(\alpha x^2),
 \qquad
 g = \Delta t - \mu x^3 c_3(\alpha x^2),
\]
donde \(x\) es una variable universal que satisface una ecuación de Kepler generalizada y
\[
 \alpha = \frac{1}{a}.
\]

Esta formulación es computacionalmente robusta y es preferida en propagadores de alta precisión que deben manejar cambios de régimen orbital o configuraciones cercanas a singularidades geométricas.

\section{Comentarios finales y perspectiva teórica}
La clase 26 cierra el ciclo de la solución analítica del problema de los dos cuerpos. Hemos pasado de la geometría estática de las cónicas a la dinámica temporal mediante la integración del teorema de áreas, la introducción de la anomalía media y la resolución de la ecuación de Kepler.

Es importante subrayar que, aunque la ecuación de Kepler es trascendental, su resolución numérica es extraordinariamente eficiente y estable. Además, la existencia de las funciones de Lagrange muestra que el espacio de estado admite una evolución lineal en términos de funciones escalares derivadas de la geometría kepleriana.

\begin{importante}
La formulación temporal del problema de dos cuerpos no solo permite predecir órbitas. También constituye la base conceptual de la teoría de perturbaciones, de la navegación espacial moderna y de los algoritmos numéricos de propagación orbital de alta precisión.
\end{importante}
\chapter{Clase 27: El problema de los tres cuerpos -- motivación y un nuevo formalismo}

El estudio de la mecánica celeste ha estado históricamente ligado a la necesidad de predecir y comprender el movimiento de los cuerpos bajo la influencia gravitatoria mutua. Desde los primeros intentos de Claudio Ptolomeo y sus epiciclos, pasando por la revolución cinemática de Johannes Kepler en los primeros años del siglo XVII, hasta la formulación dinámica de Isaac Newton en los \emph{Principia Mathematica}, la disciplina ha evolucionado desde descripciones geométricas empíricas hacia teorías físicas fundamentadas en postulados universales.

La contribución de Newton fue decisiva al establecer que el mismo conjunto de leyes que gobierna la caída de los cuerpos en la Tierra rige también el movimiento planetario. Sin embargo, la formulación newtoniana enfrenta una barrera analítica severa cuando se aplica a sistemas con más de dos cuerpos interactuando simultáneamente. En esta clase se estudia la motivación que llevó a la comunidad científica a buscar alternativas al enfoque vectorial tradicional y se presentan los fundamentos conceptuales de lo que se conoce como el \textbf{formalismo analítico o escalar de la mecánica}.

\section{El problema de los \texorpdfstring{$N$}{N} cuerpos: enfoques y limitaciones}
El problema de los \(N\) cuerpos se formula como el sistema de ecuaciones diferenciales acopladas que describe la evolución temporal de \(N\) partículas puntuales que interactúan mediante la ley de gravitación universal newtoniana:
\[
 m_i \frac{d^2 \vec{r}_i}{dt^2} = -G \sum_{\substack{j=1 \\ j \neq i}}^{N} \frac{m_i m_j}{\|\vec{r}_i - \vec{r}_j\|^3}(\vec{r}_i - \vec{r}_j),
 \quad i = 1,2,\dots,N.
\]

donde \(\vec{r}_i\) es el vector posición de la partícula \(i\), \(m_i\) su masa y \(G\) la constante de gravitación universal.

Históricamente, se han contemplado tres vías principales para abordar este problema:
\begin{enumerate}
    \item \textbf{Observación directa de la realidad:} el universo mismo actúa como un integrador numérico perfecto. Observar sistemas astronómicos reales permite seguir su evolución, pero no proporciona por sí solo una capacidad predictiva generalizable.
    \item \textbf{Solución matemática analítica:} se busca expresar las coordenadas \(\vec{r}_i(t)\) como funciones cerradas o series de la variable temporal. Para \(N \ge 3\), los resultados clásicos muestran que no existen primeras integrales algebraicas independientes adicionales a las conocidas, lo que impide una solución general cerrada.
    \item \textbf{Solución numérica:} mediante algoritmos de integración paso a paso se obtienen trayectorias aproximadas. Esta es la herramienta práctica dominante, aunque no siempre revela con claridad las simetrías o propiedades estructurales del sistema.
\end{enumerate}

La imposibilidad de una solución analítica cerrada para \(N \ge 3\) no significa que el problema sea inútil o irresoluble en la práctica. Significa que hace falta un cambio de perspectiva teórica.

\begin{importante}
El paso del problema de dos cuerpos al de tres cuerpos no es simplemente un aumento cuantitativo en el número de ecuaciones. Es un cambio cualitativo en la estructura matemática del problema.
\end{importante}

\section{Sistemas jerárquicos frente a configuraciones peculiares}
En la naturaleza, la mayoría de los sistemas gravitacionales no son configuraciones aleatorias de masas comparables. Por el contrario, tienden a organizarse en \textbf{sistemas jerárquicos}, donde las masas y las distancias relativas difieren en varios órdenes de magnitud. Ejemplos clásicos son:
\begin{itemize}
    \item el Sistema Solar, donde el Sol domina la dinámica global,
    \item los sistemas planeta--satélite,
    \item estrellas binarias con un tercer compañero distante.
\end{itemize}

En un sistema jerárquico, la dinámica puede aproximarse como la superposición de varios problemas de dos cuerpos casi independientes, complementados con correcciones perturbativas.

Sin embargo, existen configuraciones especiales o \textbf{no jerárquicas} en las que la aproximación de dos cuerpos falla de manera radical. En tales sistemas pueden aparecer comportamientos contraintuitivos:
\begin{itemize}
    \item coincidencia de períodos orbitales no explicable por una versión ingenua de la tercera ley de Kepler,
    \item distancias relativas casi constantes durante intervalos prolongados,
    \item sensibilidad extrema a las condiciones iniciales, es decir, caos determinista.
\end{itemize}

Un ejemplo moderno de este tipo de fenómenos aparece en la dinámica alrededor de los puntos de Lagrange, como \(L_2\) en el sistema Sol--Tierra, donde se ubican observatorios espaciales como el JWST. Estas configuraciones no se comprenden adecuadamente mediante una aplicación directa y simplificada de la mecánica vectorial elemental, porque las restricciones geométricas y los equilibrios dinámicos se vuelven parte esencial de la descripción.

\section{La necesidad de un nuevo formalismo: más allá de los vectores}
El formalismo vectorial de la mecánica newtoniana exige el conocimiento explícito de \textbf{todas} las fuerzas que actúan sobre cada partícula. En sistemas simples esto es directo. Pero en sistemas con restricciones mecánicas —partículas unidas por cuerdas inextensibles, cuerpos que se deslizan sobre superficies, barras rígidas, etc.— aparecen las llamadas \textbf{fuerzas de ligadura o de restricción}.

Estas fuerzas plantean dos dificultades fundamentales:
\begin{enumerate}
    \item \textbf{Son desconocidas a priori:} su magnitud y dirección dependen del estado dinámico y de las restricciones impuestas.
    \item \textbf{Complican la formulación matemática:} al incluirlas explícitamente en la segunda ley de Newton, el número de incógnitas crece y obliga a introducir ecuaciones auxiliares de ligadura que hacen el análisis más rígido y menos transparente.
\end{enumerate}

Además, el formalismo vectorial está ligado de manera natural a coordenadas cartesianas en el espacio euclidiano tridimensional. Cuando las restricciones reducen los grados de libertad del sistema, trabajar con \(3N\) coordenadas introduce una redundancia matemática innecesaria. Esta limitación motivó la búsqueda de un enfoque alternativo: un formalismo basado en cantidades escalares, capaz de eliminar automáticamente las fuerzas de restricción y de adaptarse a la geometría intrínseca del sistema.

\begin{importante}
La gran idea de la mecánica analítica no es cambiar la física, sino cambiar el lenguaje matemático con el que la física se expresa.
\end{importante}

\section{El principio de d'Alembert y el trabajo virtual}
La primera piedra angular del nuevo formalismo fue colocada por Jean le Rond d'Alembert. Su idea fundamental consistió en reemplazar el análisis directo de fuerzas por el análisis de \textbf{trabajo virtual}.

Supongamos un sistema de partículas en equilibrio estático. Si se considera un desplazamiento infinitesimal \(\delta\vec{r}_i\) compatible con las restricciones del sistema, llamado desplazamiento virtual, d'Alembert postuló que la suma del trabajo realizado por las fuerzas aplicadas es nula:
\[
 \sum_i \vec{F}_i \cdot \delta\vec{r}_i = 0.
\]

La genialidad de este principio radica en que los desplazamientos virtuales son tangentes a las superficies de restricción. Como las fuerzas de ligadura suelen ser normales a esas superficies, su producto escalar con \(\delta\vec{r}_i\) se anula automáticamente. En consecuencia, las fuerzas de restricción desaparecen de la ecuación sin necesidad de calcularlas explícitamente.

D'Alembert extendió este concepto a sistemas dinámicos reescribiendo la segunda ley de Newton en la forma
\[
 \vec{F}_i - m_i \frac{d^2\vec{r}_i}{dt^2} = \vec{0}.
\]
Interpretó el término \(-m_i\ddot{\vec{r}}_i\) como una fuerza de inercia efectiva. Así, un problema dinámico podía tratarse formalmente como si fuera un problema de equilibrio estático. El principio resultante, conocido como \textbf{principio de d'Alembert--Lagrange}, toma la forma
\[
 \sum_i \left(\vec{F}_i - m_i\frac{d^2\vec{r}_i}{dt^2}\right) \cdot \delta\vec{r}_i = 0.
\]

Esta ecuación contiene toda la información dinámica del sistema, pero aún está escrita en coordenadas cartesianas y requiere una transformación adicional para producir ecuaciones de movimiento realmente útiles.

\section{La extensión de Lagrange: fuerzas efectivas y restricciones}
Joseph-Louis Lagrange llevó la idea de d'Alembert a su máxima expresión en su \emph{Mécanique Analytique}. Si un sistema tiene \(K\) restricciones holónomas, entonces de las \(3N\) coordenadas cartesianas originales solo
\[
 M = 3N - K
\]
son verdaderamente independientes.

Lagrange introdujo entonces el concepto de \textbf{coordenadas generalizadas}
\[
 \{q_j\}_{j=1}^{M},
\]
un conjunto mínimo de variables que describe unívocamente la configuración del sistema sin redundancias. Estas coordenadas no tienen por qué ser longitudes cartesianas; pueden ser ángulos, longitudes curvilíneas u otras variables adaptadas al problema.

Al expresar los vectores posición \(\vec{r}_i\) y los desplazamientos virtuales \(\delta\vec{r}_i\) en términos de las coordenadas generalizadas \(q_j\) y sus variaciones \(\delta q_j\), el principio de d'Alembert se transforma en un sistema de ecuaciones dependiente únicamente de cantidades escalares. El resultado final, que se desarrollará con detalle en la siguiente clase, son las \textbf{ecuaciones de Euler--Lagrange}:
\[
 \frac{d}{dt}\left(\frac{\partial L}{\partial \dot{q}_j}\right) - \frac{\partial L}{\partial q_j} = 0,
 \quad j = 1,2,\dots,M,
\]
donde
\[
 L = T - V
\]
es el lagrangiano del sistema, con \(T\) la energía cinética y \(V\) la energía potencial.

La transición desde el principio de d'Alembert a las ecuaciones de Lagrange representa un cambio de paradigma: en lugar de sumar vectores fuerza en el espacio físico, se trabaja con una función escalar en el espacio de configuración. Las fuerzas de restricción no desaparecen por milagro, sino porque las coordenadas se construyen desde el principio para respetarlas.

\begin{importante}
La mecánica lagrangiana no evita las restricciones ignorándolas: las incorpora en la elección misma de las variables del problema.
\end{importante}

\section{Síntesis y transición hacia la mecánica analítica}
En esta clase se establecieron los motivos fundamentales que impulsan el abandono del formalismo vectorial tradicional en favor de un enfoque escalar:
\begin{enumerate}
    \item la complejidad analítica del problema de los \(N\) cuerpos y la inexistencia de soluciones generales cerradas,
    \item la existencia de sistemas no jerárquicos y configuraciones de equilibrio dinámico, como los puntos de Lagrange,
    \item la dificultad intrínseca para modelar fuerzas de ligadura en sistemas con restricciones,
    \item la redundancia matemática de las coordenadas cartesianas cuando los grados de libertad están reducidos.
\end{enumerate}

El principio de d'Alembert y su extensión por parte de Lagrange proporcionan un marco unificador que elimina las fuerzas de restricción por construcción geométrica y reduce la dinámica a la manipulación de funciones escalares. Este formalismo no solo simplifica la resolución de problemas clásicos, sino que además revela simetrías profundas del espacio y del tiempo y prepara el camino hacia desarrollos posteriores de la física teórica.

En la siguiente clase se aplicarán rigurosamente estos conceptos al estudio de la mecánica lagrangiana, se deducirán las ecuaciones de movimiento desde un principio variacional y se resolverán problemas concretos de mecánica celeste usando coordenadas generalizadas.
\chapter{Clase 29: El problema de los tres cuerpos y la mecánica lagrangiana}

En el desarrollo histórico y pedagógico de la mecánica celeste, el paso del problema de dos cuerpos al problema general de \(N\) cuerpos revela una brecha analítica fundamental. Mientras que el problema de dos cuerpos admite una solución cerrada basada en la conservación del momento angular, la energía y el vector de excentricidad, la introducción de una tercera partícula gravitacionalmente interactuante rompe las simetrías que permitían la integración directa. Esta clase aborda precisamente la motivación física y matemática que condujo a abandonar el enfoque puramente vectorial newtoniano en favor de un \textbf{formalismo escalar o analítico}, sentando así las bases de la mecánica lagrangiana.

El objetivo no es únicamente justificar por qué hace falta un nuevo lenguaje, sino también mostrar cómo ese nuevo lenguaje conduce naturalmente al principio de d'Alembert--Lagrange, a las coordenadas generalizadas, al lagrangiano \(L=T-U\), al principio de Hamilton y a las ecuaciones de Euler--Lagrange. La transición debe entenderse como un continuo conceptual: del problema de los tres cuerpos pasamos a la mecánica analítica porque la complejidad del primero exige la potencia estructural de la segunda.

\section{Sistemas jerárquicos y configuraciones peculiares}
Como se estableció en clases anteriores, muchos sistemas astronómicos reales pueden tratarse como \textbf{sistemas jerárquicos}, es decir, como superposiciones de subsistemas de dos cuerpos casi independientes. Esta aproximación funciona porque las masas y las escalas de distancia difieren en órdenes de magnitud tales que las perturbaciones mutuas son pequeñas y pueden incorporarse mediante teoría de perturbaciones.

Sin embargo, existen configuraciones \textbf{no jerárquicas o peculiares} donde la aproximación de dos cuerpos falla de manera drástica. Un ejemplo clásico es la dinámica cerca del punto \(L_2\) del sistema Sol--Tierra, donde se ubican observatorios como el JWST. Aunque el telescopio está más lejos del Sol que la Tierra, su período orbital coincide casi exactamente con el período terrestre. Una aplicación aislada de la tercera ley de Kepler sugeriría otra cosa, pero la presencia simultánea de dos potenciales gravitacionales modifica la dinámica efectiva del sistema.

Otro ejemplo ilustrativo es un sistema de tres masas dispuestas con condiciones iniciales cuidadosamente escogidas. Al integrar sus ecuaciones de movimiento, pueden observarse períodos casi coincidentes, equilibrios dinámicos precarios y una sensibilidad extrema a pequeñas variaciones iniciales, es decir, \textbf{caos determinista}. Estos comportamientos muestran que la suma vectorial directa de fuerzas newtonianas, aunque correcta en principio, deja de ser una herramienta analítica cómoda y transparente.

\begin{importante}
El problema de los tres cuerpos no exige cambiar las leyes físicas; exige cambiar el lenguaje matemático con el que describimos esas leyes.
\end{importante}

\section{Limitaciones del formalismo vectorial}
La mecánica newtoniana clásica se basa en el postulado
\[
\vec{F}_{\text{total}} = \frac{d\vec{p}}{dt} = m\vec{a}.
\]
Para predecir la evolución de un sistema, debemos conocer explícitamente \textbf{todas} las fuerzas que actúan sobre cada partícula. En sistemas simples esto no representa mayor dificultad. Pero cuando aparecen \textbf{restricciones mecánicas} —cuerdas inextensibles, superficies rígidas, barras, péndulos, etc.— surgen las llamadas \textbf{fuerzas de ligadura o de restricción}.

Estas fuerzas presentan dos dificultades estructurales:
\begin{enumerate}
    \item \textbf{Son desconocidas a priori:} su magnitud y dirección dependen del estado instantáneo del sistema y de la geometría de la restricción.
    \item \textbf{Introducen redundancia matemática:} si un sistema tiene \(N\) partículas, se requieren \(3N\) coordenadas cartesianas. Pero si existen \(K\) restricciones holónomas, solo
\[
M = 3N - K
\]
coordenadas son realmente independientes.
\end{enumerate}

Trabajar entonces con \(3N\) ecuaciones vectoriales acopladas, incluyendo fuerzas desconocidas de restricción, vuelve innecesariamente engorroso un problema cuya geometría real puede tener muchos menos grados de libertad.

\section{El principio de los trabajos virtuales y la idea de d'Alembert}
La primera piedra del nuevo formalismo fue colocada por Jean le Rond d'Alembert. Su contribución fundamental consistió en reemplazar el análisis directo de fuerzas por el análisis de \textbf{trabajo virtual}.

\subsection{Desplazamiento virtual}
Un \textbf{desplazamiento virtual} \(\delta\vec{r}_i\) es un desplazamiento infinitesimal imaginario de la partícula \(i\)-ésima que:
\begin{itemize}
    \item se realiza instantáneamente, de modo que \(\delta t = 0\),
    \item es compatible con las restricciones del sistema en el instante considerado,
    \item no implica evolución temporal real, sino una variación geométrica admisible.
\end{itemize}

\subsection{Principio de los trabajos virtuales}
Para un sistema en equilibrio estático, d'Alembert postuló que la suma del trabajo realizado por las fuerzas aplicadas sobre todos los desplazamientos virtuales compatibles es nula:
\[
\sum_i \vec{F}_i \cdot \delta\vec{r}_i = 0.
\]

La potencia de este principio radica en que los desplazamientos virtuales son tangentes a las superficies de restricción. Como las fuerzas de ligadura suelen ser normales a esas superficies, su producto escalar con \(\delta\vec{r}_i\) se anula. Así, las fuerzas de restricción desaparecen automáticamente de la formulación.

\subsection{Extensión a la dinámica: principio de d'Alembert--Lagrange}
D'Alembert extendió esta idea a sistemas en movimiento reescribiendo la segunda ley de Newton como
\[
\vec{F}_i - m_i\ddot{\vec{r}}_i = \vec{0}.
\]
Interpretando el término \(-m_i\ddot{\vec{r}}_i\) como una fuerza de inercia efectiva, convirtió el problema dinámico en una especie de problema de equilibrio virtual. El resultado es el \textbf{principio de d'Alembert--Lagrange}:
\[
\sum_i \left(\vec{F}_i - m_i\ddot{\vec{r}}_i\right)\cdot\delta\vec{r}_i = 0.
\]

Esta ecuación ya contiene toda la dinámica del sistema, pero todavía está escrita en coordenadas cartesianas. El siguiente paso consiste en expresar esos desplazamientos virtuales en variables intrínsecas adaptadas a la geometría real del problema.

\section{Restricciones holónomas, grados de libertad y variables generalizadas}
Una restricción puede expresarse matemáticamente como
\[
f_k(\vec{r}_1,\vec{r}_2,\dots,\vec{r}_N,t)=0, \qquad k=1,2,\dots,K.
\]
Cuando una restricción puede escribirse como una igualdad algebraica entre coordenadas y tiempo, se dice que es \textbf{holónoma}. Para estos sistemas, el número de coordenadas independientes es precisamente
\[
M = 3N-K,
\]
que coincide con el número de \textbf{grados de libertad}.

En lugar de trabajar con las \(3N\) coordenadas cartesianas redundantes, introducimos un conjunto de \(M\) variables independientes
\[
\{q_j\}_{j=1}^{M},
\]
llamadas \textbf{coordenadas generalizadas}. Estas variables pueden ser longitudes, ángulos u otras cantidades físicas capaces de parametrizar el espacio de configuración sin redundancias.

Las posiciones cartesianas pasan entonces a escribirse como funciones de estas coordenadas:
\[
\vec{r}_i = \vec{r}_i(q_1,q_2,\dots,q_M,t).
\]
La condición fundamental es que las \(q_j\) sean independientes, de modo que las variaciones virtuales \(\delta q_j\) puedan considerarse libres entre sí.

\subsection{Transformación de desplazamientos virtuales}
Aplicando la regla de la cadena y recordando que \(\delta t=0\), se obtiene:
\[
\delta\vec{r}_i = \sum_{j=1}^{M} \frac{\partial\vec{r}_i}{\partial q_j}\,\delta q_j.
\]
Esta relación conecta el espacio físico cartesiano con el espacio abstracto de configuración. Al sustituirla en el principio de d'Alembert--Lagrange, la dinámica queda proyectada sobre las direcciones permitidas por las variables generalizadas.

\begin{importante}
Las coordenadas generalizadas no son una simple comodidad algebraica. Son la expresión matemática de los verdaderos grados de libertad del sistema.
\end{importante}

\section{Hacia las ecuaciones de Euler--Lagrange}
Sustituyendo la expresión de \(\delta\vec{r}_i\) en el principio de d'Alembert--Lagrange, obtenemos
\[
\sum_{j=1}^{M}\left[\sum_{i=1}^{N}\vec{F}_i\cdot\frac{\partial\vec{r}_i}{\partial q_j} - \sum_{i=1}^{N}m_i\ddot{\vec{r}}_i\cdot\frac{\partial\vec{r}_i}{\partial q_j}\right]\delta q_j = 0.
\]

El primer término define la \textbf{fuerza generalizada}
\[
Q_j \equiv \sum_{i=1}^{N} \vec{F}_i\cdot\frac{\partial\vec{r}_i}{\partial q_j}.
\]
El segundo término, usando identidades de cálculo y la definición de energía cinética total
\[
T = \frac{1}{2}\sum_{i=1}^{N} m_i\dot{\vec{r}}_i^{\,2},
\]
puede transformarse en
\[
\sum_{i=1}^{N}m_i\ddot{\vec{r}}_i\cdot\frac{\partial\vec{r}_i}{\partial q_j}
=
\frac{d}{dt}\left(\frac{\partial T}{\partial \dot{q}_j}\right)-\frac{\partial T}{\partial q_j}.
\]

Con esto, y aprovechando la independencia de las variaciones \(\delta q_j\), se llega a la forma general
\[
\frac{d}{dt}\left(\frac{\partial T}{\partial \dot{q}_j}\right)-\frac{\partial T}{\partial q_j}=Q_j,
\qquad j=1,2,\dots,M.
\]

Si además las fuerzas son conservativas, existe un potencial \(U\) tal que
\[
Q_j = -\frac{\partial U}{\partial q_j}.
\]
Entonces resulta natural definir el \textbf{lagrangiano}
\[
L(q_j,\dot{q}_j,t)=T-U,
\]
y las ecuaciones de movimiento adoptan la forma canónica
\[
\frac{d}{dt}\left(\frac{\partial L}{\partial \dot{q}_j}\right)-\frac{\partial L}{\partial q_j}=0,
\qquad j=1,2,\dots,M.
\]

Estas son las \textbf{ecuaciones de Euler--Lagrange}. A partir de este punto, la dinámica queda enteramente codificada en una función escalar.

\section{El principio de Hamilton y la estructura variacional}
La formulación lagrangiana puede reinterpretarse de manera aún más profunda mediante el \textbf{principio de Hamilton}. Dado un lagrangiano \(L\), definimos la acción
\[
S[q_j(t)] = \int_{t_1}^{t_2} L(q_j,\dot{q}_j,t)\,dt.
\]

El principio de Hamilton afirma que la trayectoria real seguida por el sistema es aquella para la cual la acción es estacionaria frente a variaciones compatibles de la trayectoria, manteniendo fijos los extremos:
\[
\delta S = \delta\int_{t_1}^{t_2}L\,dt = 0.
\]

Aplicando el cálculo de variaciones a esta condición se recuperan exactamente las ecuaciones de Euler--Lagrange. La dinámica deja entonces de interpretarse como el resultado de una suma vectorial instantánea de fuerzas y pasa a verse como un problema variacional global en el espacio de configuraciones.

\begin{importante}
La mecánica lagrangiana muestra que la naturaleza puede describirse no solo por ecuaciones diferenciales, sino también por principios de extremación. Esta es una de las ideas más profundas de toda la física teórica.
\end{importante}

\section{Simetrías, coordenadas cíclicas y leyes de conservación}
Una de las grandes virtudes del formalismo lagrangiano es que revela constantes de movimiento a partir de simetrías, sin necesidad de resolver explícitamente las ecuaciones diferenciales.

Si el lagrangiano no depende de una coordenada generalizada \(q_k\), es decir,
\[
\frac{\partial L}{\partial q_k}=0,
\]
entonces \(q_k\) es una \textbf{coordenada cíclica} y la ecuación correspondiente se reduce a
\[
\frac{d}{dt}\left(\frac{\partial L}{\partial \dot{q}_k}\right)=0.
\]
Esto implica la conservación del \textbf{momento generalizado conjugado}
\[
p_k \equiv \frac{\partial L}{\partial \dot{q}_k} = \text{constante}.
\]

De manera similar, si el lagrangiano no depende explícitamente del tiempo, se conserva la cantidad
\[
h = \sum_{j=1}^{M}\dot{q}_j\frac{\partial L}{\partial \dot{q}_j} - L,
\]
que, bajo condiciones usuales, coincide con la energía mecánica total.

Estos resultados son manifestaciones particulares del \textbf{teorema de Noether}: a cada simetría continua del lagrangiano le corresponde una cantidad conservada. Translaciones espaciales se asocian al momento lineal, rotaciones al momento angular y traslaciones temporales a la energía.

\section{Aplicación a la mecánica celeste}
El formalismo lagrangiano no es solo una reformulación elegante; es una herramienta conceptual y computacional superior para la mecánica celeste.

\subsection{Lagrangiano del problema de \texorpdfstring{$N$}{N} cuerpos}
Para un sistema de \(N\) partículas que interactúan exclusivamente por gravedad newtoniana, no existen restricciones holónomas adicionales. Las coordenadas cartesianas \(\{\vec{r}_i\}\) pueden usarse como variables generalizadas. El lagrangiano se escribe entonces como
\[
L = \sum_{i=1}^{N}\frac{1}{2}m_i\dot{\vec{r}}_i^{\,2} + \frac{1}{2}\sum_{i=1}^{N}\sum_{\substack{j=1 \\ j \neq i}}^{N}\frac{Gm_im_j}{\|\vec{r}_i-\vec{r}_j\|}.
\]

El factor \(\tfrac{1}{2}\) evita la doble contabilidad de pares de interacción. Aplicando las ecuaciones de Euler--Lagrange, se recuperan las ecuaciones de movimiento newtonianas, pero ahora insertas en un marco donde las simetrías globales del sistema son más fáciles de identificar.

\subsection{Reducción al problema de dos cuerpos y potencial efectivo}
Cuando se aísla un par dentro de un sistema jerárquico, puede introducirse la coordenada del centro de masa \(\vec{R}_{\text{CM}}\) y el vector relativo \(\vec{r}=\vec{r}_1-\vec{r}_2\). La energía cinética se descompone de modo natural y, tras eliminar la coordenada cíclica asociada al centro de masa, se obtiene un lagrangiano efectivo para el movimiento relativo:
\[
L_{\text{rel}} = \frac{1}{2}\mu\dot{\vec{r}}^{\,2} + \frac{\mu\mathcal{G}}{r},
\qquad
\mathcal{G}=G(m_1+m_2).
\]

Al pasar luego a coordenadas polares o esféricas y fijar el plano orbital, aparece el \textbf{potencial efectivo}
\[
V_{\text{eff}}(r) = -\frac{\mu\mathcal{G}}{r} + \frac{h^2}{2r^2},
\]
que permite analizar de manera cualitativa la existencia de órbitas ligadas, no ligadas, periapsis, apoapsis y órbitas circulares estables, todo ello sin integrar directamente las ecuaciones de movimiento.

\subsection{Perturbaciones y elementos osculantes}
Cuando un tercer cuerpo o una fuerza no newtoniana introduce una perturbación pequeña, el lagrangiano se modifica ligeramente. A partir de él puede desarrollarse de manera sistemática la teoría de perturbaciones lagrangiana y derivarse la evolución temporal de los elementos orbitales osculantes. De este modo, la mecánica lagrangiana se convierte en el marco natural para pasar del problema kepleriano idealizado a la dinámica celeste realista.

\section{Síntesis final}
La transición del problema de los tres cuerpos a la mecánica lagrangiana no es un cambio de tema, sino un cambio de nivel conceptual. En esta clase se ha visto que:
\begin{enumerate}
    \item la complejidad del problema de \(N\) cuerpos obliga a abandonar la dependencia explícita en fuerzas de restricción y coordenadas redundantes,
    \item el principio de d'Alembert--Lagrange reformula la dinámica como una condición escalar sobre desplazamientos virtuales compatibles,
    \item las coordenadas generalizadas describen el sistema con el número mínimo de variables independientes,
    \item el lagrangiano \(L=T-U\) unifica cinemática y dinámica en una sola función escalar,
    \item el principio de Hamilton y el teorema de Noether revelan la estructura variacional y simétrica profunda de la mecánica,
    \item en mecánica celeste, este formalismo permite describir desde el problema de dos cuerpos hasta la teoría de perturbaciones con una claridad muy superior al enfoque vectorial elemental.
\end{enumerate}

En consecuencia, la mecánica lagrangiana no debe verse únicamente como una técnica alternativa. Es el lenguaje natural para estudiar sistemas gravitacionales complejos y constituye el puente indispensable hacia formulaciones aún más generales, como la mecánica hamiltoniana, la teoría de perturbaciones moderna y, en última instancia, la física del siglo XX y XXI.

% ---------- FINAL ----------
\backmatter
% Puedes agregar bibliografía o anexos aquí si deseas.

\end{document}